% paper.tex -- outline of the revised paper (BIOINF-2026-1434 revision). Outline phase: locked information only.
% Build:  tectonic paper.tex   (bibliography: refs.bib, natbib plainnat). ASCII-only source.
% Built from branch final-review-fixes at commit 08f1e3c (2026-10-14); see 'How to read this outline'.

\documentclass[11pt]{article}
\usepackage[margin=2.2cm]{geometry}
\usepackage{amsmath,amssymb}
\usepackage{booktabs,tabularx,longtable,array}
\usepackage{enumitem}
\usepackage{needspace,etoolbox}
\pretocmd{\paragraph}{\needspace{4\baselineskip}}{}{}  % keep run-in headings with their first lines
\usepackage[dvipsnames]{xcolor}
\usepackage[authoryear,round]{natbib}
\usepackage{url}
\usepackage{hyperref}
\hypersetup{colorlinks=true,linkcolor=black,citecolor=MidnightBlue,urlcolor=MidnightBlue}
\setlist{nosep,leftmargin=1.4em}
\setlength{\emergencystretch}{3em}
\newcolumntype{L}{>{\raggedright\arraybackslash}X}
\DeclareUrlCommand\file{\urlstyle{tt}}
% Outline status macros
%   plain text + \src{...}   = LOCKED (user decision in CONSTRAINTS.md, rule frozen at tag prereg-tier12-v2, or verified fact)
%   \pending{what}{when}      = not locked; 'when' names the event that will lock it
%   \unverified{key}          = citation lead whose claim has not been checked against the source
\newcommand{\src}[1]{{\scriptsize\textcolor{gray}{[#1]}}}
\newcommand{\pending}[2]{{\color{BurntOrange}\textbf{[PENDING:} #1 \textit{--- locks when: #2}\textbf{]}}}
\newcommand{\unverified}[1]{{\color{red}[?#1]}}


\begin{document}
\begin{center}
{\LARGE\bfseries Paper outline}\\[4pt]
{\large Wasserstein critic versus discriminator in adversarial batch correction on scVI}\\[6pt]
\pending{title}{your decision}\quad \pending{authors and target journal}{your decision}\\[2pt]
{\small Outline of the revision of Bioinformatics submission BIOINF-2026-1434; repository state: commit 08f1e3c, 2026-10-14.}
\end{center}

\section*{How to read this outline}
\begin{itemize}
\item Three status forms. Plain text followed by a grey bracket \src{source} is \textbf{locked}: a user decision in the decision ledger (\file{CONSTRAINTS.md}), a rule frozen at tag prereg-tier12-v2 (commit c88cce3), or a fact verified in a file. \pending{what is still open}{the event that will lock it} marks text that is \textbf{not locked}. \unverified{key} marks a citation key whose claim has not yet been checked against its source; none remains. Author-year citations are checked: each claim against the source text, recorded in \file{paper/evidence.md}, which names the citations checked against the abstract or metadata only and those read in a preprint version.
\item SI-nn: entry nn of the decision ledger \file{CONSTRAINTS.md} (77 entries). Each quotes the user's answer, except SI-42, a lead-agent decision that is not locked, and the sentence of SI-17 marked as an agent-derived corollary. SI-35 and SI-59 are retired (replaced by SI-46 and SI-61).
\item PREREG \S n: section n of \file{docs/PREREG.md}, the pre-registered decision rules. Frozen at tag prereg-tier12-v2 (commit c88cce3) and byte-identical at 08f1e3c: \S0--\S2, \S4, \S6--\S8, the Disclosure paragraph, and 11 of the 12 section 9 entries dated 2026-10-03/04 \src{\file{docs/PREREG.md} at c88cce3 and 08f1e3c, compared}. At 08f1e3c, \S3, \S5 and the 2026-10-03 entry also carry bracketed notes dated 2026-10-13/14 (version 3): the \S3 note and the 2026-10-03 correction implement SI-56, SI-66 and SI-53; the \S5 note (Holm over six results) is an agent proposal. The section 9 entries dated 2026-10-13 and 2026-10-14 are version 3, which is not frozen.
\item ANALYSIS\_SPEC \S n: section n of \file{docs/ANALYSIS_SPEC.md}, the draft specification of the X1 analyses. It holds user decisions (SI-68 to SI-77) and agent proposals that wait for the user's review of the full diff.
\item PAPER\_PLAN \S n: section n of \file{docs/PAPER_PLAN.md}, the design-time paper plan of 2026-10-02, partly superseded.
\item NB-yyyymmdd-nn: entry nn of the lab notebook page \file{notebook/yyyy-mm-dd.md}.
\item I0nn: finding nn of the final pre-publication review of 2026-10-13 (\file{docs/review_final/REVIEW_FINAL.md}; full list in \file{D/review/lead/verified_issues.csv}). D-nn: finding nn of that review's recheck of the design documents (\file{D/review/rereview/design_fable/REVIEW_design_recheck.md}).
\item file:line: a line of a repository file at commit 08f1e3c (branch final-review-fixes, 2026-10-14). Files in \file{scripts/} are named without the folder; \file{wcd/} stands for \file{src/wcd_vae/wcd/}. \file{D/} is the project data folder.
\item diff review: the user's review of the full diff of branch final-review-fixes before the tag prereg-tier12-v3 is made \src{SI-52; SI-60}.
\end{itemize}
\subsection*{Glossary}
{\small
\begin{longtable}{>{\raggedright\arraybackslash}p{0.17\linewidth}>{\raggedright\arraybackslash}p{0.76\linewidth}}
\toprule
Label & Plain meaning \\
\midrule
\endhead
A1 & Pilot of penalty strengths: atac\_small, immune and sim1, both decoders, seeds 100--102; plain scVI plus the six adversarial arms at ten strengths (0.1 to 3000). Its scores set the X1 strengths (rule R1). \src{PREREG \S0, \S2; SI-16} \\
A2 & Pilot of the adversary input: a posterior sample instead of the posterior mean (atac\_small and immune; discriminator, reference and pooled critics; both decoders; seeds 100--102). \src{PREREG \S3; SI-27} \\
A3 & Pilot of per-dimension standardisation of the adversary input (same tasks, decoders and seeds; discriminator, reference, pooled, MMD, Sinkhorn). \src{PREREG \S3; SI-27} \\
X1 & Core benchmark: 8 tasks, both decoders, seeds 0--4; plain scVI, the six adversarial arms at six strengths each, and scvi-tools' own adversary (conditioned decoder only); 3,000 fits. \src{manifest; SI-16} \\
X3 & Composition shift: one batch of a task loses two (sim2: one) chosen cell types at doses 0, 50, 80, 95 and 100\%, with an importance-weighted control. \src{SI-25, SI-41; manifest} \\
X6 & Regularisation and update budget: the discriminator plain, spectrally normalised or with an R1 gradient penalty (Mescheder et al.; not rule R1); the pooled critic at 1, 5 and 10 critic steps, and spectrally normalised at 1 and 5 steps. \src{manifest; SI-22} \\
X7 & Reference choice, as designed in the manifest of record: every batch in turn as the reference (atac\_small, pancreas, sim1), for the reference critic, a fixed-reference critic, reference-anchored MMD and reference-anchored JS \src{manifest; SI-23}; \pending{redesign (NO-GO as designed: D-06, I089)}{during X1, before any X1 result is examined (SI-68)} \\
X8 & Batch-count scaling, as designed in the manifest of record: V = 2, 4, 8 or 16 batches at a fixed total number of cells per task (immune\_hum\_mou 16,000, lung 6,000, sim2 3,000) \src{manifest}; \pending{redesign (NO-GO as designed: D-07, I015, I077, I078)}{during X1, before any X1 result is examined (SI-68)} \\
X10 & Downstream biology at matched batch removal; uses X1 latents and models, no fits of its own; exploratory. \src{SI-68; ANALYSIS\_SPEC \S0} \\
X12 & Architecture and hyperparameter robustness: an 8-run fractional factorial over latent size, adversary width, adversary learning rate, minibatch size and decoder (scVI's default or LinearSCVI's linear decoder). \src{\file{build_paper_manifest.py}:259; \file{scvi_adversarial_plan.py}:516--519; manifest} \\
X13 & Standard integrators: Harmony, Scanorama and PCA (run once per setting), scANVI and sysVI (seeds 0--4). \src{SI-34, SI-36, SI-46; manifest} \\
X15 & KL warm-up sensitivity: scvi-tools' 400-epoch KL warm-up against a warm-up that completes within the task's epochs. \src{SI-44} \\
R0 to R5 (results sections) & R0 premise; R1 integration versus conservation; R2 scalability with batch count; R3 sensitivity to penalty strength and design factors; R4 reference selection; R5 placement against standard integrators. \src{PAPER\_PLAN \S3 table} \\
rules R1 to R4 & The pre-registered decision rules. Rule R1 sets each family's six X1 strengths from A1; rule R2 the adversary input; rule R3 standardisation; rule R4 the matched strengths. This outline always writes ``rule R1'' to keep it apart from results section R1. \src{PREREG \S2--4} \\
P1 to P4 & Primary endpoints. P1: biology score at the matched batch score b*, reference critic minus discriminator, one result per decoder, with its trade-off area \src{SI-68; SI-70; SI-72}. P2: how often runs collapse \src{SI-75}, with the X3 dose-response specified later \src{SI-68}. P3 (batch-count scaling, X8) and P4 (reference choice, X7): \pending{definitions (PAPER\_PLAN \S3 proposes the change of the critic-minus-discriminator difference with batch count, and a formulation-by-coverage interaction)}{X7/X8 redesign, during X1 (SI-68)} \\
B & Batch score: the unweighted mean of five scIB batch metrics, unscaled; higher means more batch effect removed. \src{PREREG \S1} \\
C & Biology score: the unweighted mean of the task's scIB biology metrics, unscaled; higher means more biology kept. \src{PREREG \S1} \\
b0, c0 & Mean B and mean C of plain scVI (penalty strength 0) in one task with one decoder. \src{PREREG \S4; ANALYSIS\_SPEC (words used)} \\
b\_common & The batch score that every adversarial arm reaches: over the six arms, the lowest of each arm's highest failure-free seed-mean B. \src{PREREG \S4; SI-21} \\
b* & Matched batch-removal target: b0 + 0.5 (b\_common $-$ b0), half-way from plain scVI to the batch removal every arm reaches. \src{SI-21} \\
low-signal block & A task with one decoder where b\_common $-$ b0 $<$ 0.02 (flag bstar\_low\_signal), so that b* cannot be placed reliably; not evaluable for P1 or rules R2/R3, reported separately. \src{SI-56; SI-66; PREREG \S4} \\
collapse & In rule R1: a strength at which a seed failed or the seed-mean C fell by 0.10 or more against plain scVI. In P2: a run that failed, or whose C is 0.10 or more below the same seed's plain-scVI run. \src{PREREG \S2; SI-75} \\
run (fit), failed & One model trained with one seed. Failed: the training loss became non-finite (diverged) or the saved latent has a non-finite value (nonfinite\_latent); infrastructure errors are rerun, not failures. \src{PREREG \S1} \\
point & One penalty strength of one method in one task with one decoder; in X1, 5 runs. \src{ANALYSIS\_SPEC (words used)} \\
cell & In the X1 analyses: one task with one decoder. In rules R2/R3: one task, decoder and arm. \src{ANALYSIS\_SPEC (words used); SI-27} \\
decoder & Conditioned: the decoder receives the batch label, as scVI is usually used. Unconditioned: the decoder gets no batch label, so the adversary is the only batch correction. \src{\file{scvi_adversarial_plan.py}:534, :537; \file{tests/scvi/test_adversarial_plan_v2.py}:118} \\
arm & One method in the comparison: none (plain scVI, strength 0); discriminator (JS-type, after the Jensen--Shannon (JS) form of the objective of the original generative adversarial network (GAN)); reference, pooled and barycenter critics (W1, Wasserstein-1); mmd (maximum mean discrepancy); sinkhorn (Sinkhorn divergence); scvi\_adv (scvi-tools' own adversary). \src{manifest; \file{scvi_adversarial_plan.py}:4--17} \\
variant arms & Follow-up versions of an arm: discriminator\_sn and pooled\_sn (spectrally normalised), discriminator\_r1 (R1 gradient penalty of Mescheder et al., unrelated to rule R1), reference\_fixed (reference cells held fixed), discriminator\_ref and mmd\_ref (reference-anchored JS and MMD). \src{manifest; \file{scvi_adversarial_plan.py}:12--24; SI-22; SI-23} \\
penalty strength, lambda & The weight $\lambda$ of the alignment term that an arm adds to the scVI loss; $\lambda = 0$ is plain scVI, except in scvi\_adv rows (manifest $\lambda$ 0), whose classifier loss is scaled by $1 -$ KL weight instead. \src{\file{scvi_adversarial_plan.py}:436, :177; manifest} \\
family & Arms that share one strength grid: JS \{discriminator\}; W1 \{reference, pooled, barycenter\}; MMD \{mmd\}; Sinkhorn \{sinkhorn\}; one grid per decoder (8 grids). \src{SI-18} \\
task, real task, dataset family & Task: one of the eight prepared scIB datasets. Real tasks: pancreas, lung, immune, immune\_hum\_mou, atac\_small, atac\_large; sim1 and sim2 are simulations. Dataset families of real tasks: pancreas; lung; immune (immune, immune\_hum\_mou); ATAC (atac\_small, atac\_large). \src{SI-02; SI-73; \file{prereg_rules.py}:66--68} \\
seeds & A1--A3: 100--102; X1 and the neural X13 baselines (scANVI, sysVI): 0--4, Harmony, Scanorama and PCA once at seed 0; follow-ups X3, X6, X7, X8, X12, X15: 10--12. The three sets are disjoint. \src{SI-26; SI-44; SI-34; SI-46; SI-36; manifest} \\
CR-nn & A finding of the code review of tag prereg-tier12-v1 (2026-10-03); findings of the re-review of its fixes (2026-10-04) are numbered RR-nn. \src{\file{docs/PREREG.md}:167} \\
KL warm-up & KL: Kullback--Leibler. The weight of the KL term of the scVI loss rises linearly from 0 over 400 epochs (scvi-tools default) and ends at $(E-1)/400$ for a task with $E \leq 400$ epochs: 0.998 at 400 epochs, 0.20--0.61 on the four shorter tasks [S-E3]. \src{\file{build_paper_manifest.py}:104; \file{fit_paper_config.py}:123} \\
Tier 1, Tier 2 & The design-time scope tiers; Tier 2 adds X3, X6 and X12. \src{SI-05; PAPER\_PLAN \S4} \\
metric names & PCR: principal-component regression; ASW: average silhouette width; iLISI, cLISI: integration and cell-type local inverse Simpson's index; kBET: k-nearest-neighbour batch-effect test; NMI: normalised mutual information; ARI: adjusted Rand index; F1: isolated-label F1 score. \src{\file{score_scib_native.py}:51--52} \\
other abbreviations & GP: gradient penalty; WGAN-GP: Wasserstein GAN (generative adversarial network) with gradient penalty; POT: Python Optimal Transport; OT: optimal transport; HVG: highly variable gene; SD: standard deviation; JHPCE: the remote computing cluster with the L40S GPUs. \src{SI-39 (JHPCE)} \\
\bottomrule
\end{longtable}
}


\tableofcontents
\clearpage

\section*{Abstract (outline)}
\addcontentsline{toc}{section}{Abstract (outline)}
\begin{itemize}
\item Question: does replacing the Jensen--Shannon (JS) discriminator of adversarial batch correction with a Wasserstein-1 (W1) critic change the trade-off between batch removal and biology conservation on scVI, with maximum mean discrepancy (MMD) and Sinkhorn as non-adversarial comparators? \src{SI-09; SI-12; SI-18; manifest}
\item Design: six adversarial arms (discriminator; reference, pooled and barycenter W1 critics; MMD; Sinkhorn) and scvi-tools' own adversary on one scvi-tools default scVI backbone \src{SI-10; manifest}; batch-conditioned and unconditioned decoders \src{SI-11; SI-16}; the eight scIB tasks, fitted as distributed \src{SI-02; SI-03; SI-15}; scIB metrics without min--max scaling \src{SI-08; PREREG \S1}.
\item Fixed in advance: penalty-strength grids from a pilot on seeds disjoint from the main experiment, and decision rules R1--R4 \src{PREREG \S0--\S4 and Disclosure; SI-13; SI-16}; the main-experiment analyses, locked before any of its results is examined \src{SI-52; SI-68}; \pending{freeze of that analysis plan}{diff review, then tag prereg-tier12-v3}.
\item Primary comparison (P1): biology score at matched batch removal, reference critic minus discriminator, one result per decoder \src{SI-21; SI-68; SI-70}.
\item Main result: \pending{main result}{X1 scored and the single analysis run}
\item Implication: \pending{guidance for scVI users, including the case that no adversary improves on plain scVI}{X1 and X13 analysed}
\item Availability: code in the GitHub repository kreid415/wasserstein-critic-deconfounding \src{git remote}; \pending{tag and commit of record}{tag prereg-tier12-v3}; scIB task files with checksums [S-D2].
\end{itemize}

\section{Introduction}\label{sec:intro}

\paragraph{Batch effects.}
\begin{itemize}
\item Batch effects are technical variation from laboratory conditions, reagent lots, personnel, instruments, time or pipeline; they mislead when correlated with the outcome \citep{leek2010tackling,yu2024assessing}.
\item In single-cell RNA sequencing (scRNA-seq), separately processed replicates differ substantially \citep{tung2017batch}; batch effects are expected to remain relevant, and scRNA-seq brings new challenges for modelling them \citep{goh2022batch}.
\item Correction can fail both ways: uncorrected batch effects can mislead, and over-correction can hinder biomedical discovery \citep{yu2024assessing}.
\item Integration is scored on two axes: batch removal and conservation of biological variation \citep{luecken2022benchmarking}.
\end{itemize}

\paragraph{Adversarial integration.}
\begin{itemize}
\item An encoder is trained against an adversary that predicts the batch, or another confounder such as age, from the latent code \citep{ganin2016domain,dincer2020adversarial}.
\item A discriminator targets the Jensen--Shannon (JS) divergence: for the optimal discriminator, the generative adversarial network (GAN) objective equals $-\log 4$ plus twice the JS divergence between the data and generator distributions \citep[Sec.~4.1, Eq.~6, in the proof of Thm.~1]{goodfellow2014generative}.
\item A Wasserstein-1 (W1) critic estimates the earth mover's distance through its 1-Lipschitz dual \citep{arjovsky2017wasserstein}; WGAN-GP (Wasserstein GAN with gradient penalty) enforces the constraint softly \citep[Sec.~4]{gulrajani2017improved}.
\end{itemize}

\paragraph{JS versus W1, stated precisely.}
\begin{itemize}
\item Reading the theory for batch correction: the encoder plays the generator, the latent codes of two batches are the two distributions to be aligned, and the discriminator or critic is the adversary \src{\file{scvi_adversarial_plan.py}:346--360, :433--439}.
\item W1 is continuous for continuous generators, and differentiable almost everywhere for locally Lipschitz ones (with a regularity assumption); for the JS and Kullback--Leibler (KL) divergences these statements fail in general \citep[Thm.~1]{arjovsky2017wasserstein}.
\item On disjoint parallel lines JS $=\log 2$ at every nonzero offset, while W1 equals the offset \citep[Example~1]{arjovsky2017wasserstein}.
\item For disjoint compact supports (Thm.~2.1), or supports on closed manifolds that do not perfectly align and do not have full dimension (Thm.~2.2), a perfect discriminator exists whose gradient is zero on the supports (everywhere in Thm.~2.1, at almost every point in Thm.~2.2) \citep[Thms.~2.1--2.2]{arjovsky2017towards}.
\item The saturating loss $\log(1-D)$ then has vanishing generator gradients near the optimal discriminator \citep[Thm.~2.4]{arjovsky2017towards}.
\item The non-saturating loss $-\log D$ need not have vanishing gradients: for distributions with densities, at the optimal discriminator its expected gradient is that of $\mathrm{KL}(P_g\Vert P_r)-2\,\mathrm{JS}(P_g\Vert P_r)$ (Thm.~2.5); if the discriminator is the optimal one plus independent Gaussian noise at every point (and in its gradient), each gradient coordinate is Cauchy-distributed, with infinite mean and variance (Thm.~2.6, under the conditions of Thms.~2.1--2.2; the authors call this independence assumption too strong) \citep[Sec.~2.2.2]{arjovsky2017towards}.
\item Continuous noise added to the discriminator input gives both distributions a density (Thm.~3.1); with Gaussian noise the densities are positive everywhere, so the supports overlap (Cor.~3.1) \citep{arjovsky2017towards}.
\item scVI (single-cell variational inference) encodes each cell as a Gaussian in the latent space, with a mean and a variance; the posterior mean is that mean, and a posterior sample adds Gaussian noise with that variance; the noise scale is therefore cell-specific, a case Thm.~3.1 does not cover directly \src{scvi-tools 1.4.2: scvi/nn/\_base\_components.py:304--305}.
\item In the main experiment X1, scvi-tools' own adversary (arm scvi\_adv) is given posterior samples \src{scvi-tools 1.4.2: scvi/module/\_vae.py:390, :413, scvi/train/\_trainingplans.py:689; \file{scvi_adversarial_plan.py}:348--351}; the manifest's adv\_input column reads `mean' for its 40 X1 rows, a value this arm does not read \src{\file{manifests/paper_manifest_stock_pilot_u5_b10_v3.tsv}}.
\item The six arms implemented in this project (discriminator; reference, pooled and barycenter W1 critics; MMD; Sinkhorn) are given the posterior mean by default \src{SI-19; \file{scvi_adversarial_plan.py}:359}. X1 keeps this default and keeps standardisation off \src{SI-56; SI-66}: the pilot rules that could switch the adversary input to posterior samples (rule R2) or standardise it per latent dimension (rule R3) keep their defaults while a block they use is flagged low-signal, that is, while the batch removal that all six arms reach exceeds plain scVI's by less than 0.02. \pending{formal re-run of rules R2 and R3}{after the Sinkhorn redo (SI-57)}
\item The frozen record of the strength-matching rule (rule R4), run on the penalty-strength pilot (A1), flags two blocks, both with the conditioned decoder: immune (difference 0.0188) and sim1 (0.0069). Rules R2 and R3 are run on pilots A2 and A3, which cover atac\_small and immune only, so only immune bears on them \src{\file{docs/prereg/r4_a1_record.json}; manifest}.
\item Refitting Sinkhorn (the Sinkhorn arm was corrected after the pilot, so its pilot fits are redone) cannot clear either flag: the common level is set by the arm that reaches the least batch removal (barycenter in five of the six blocks, MMD in sim1 with the conditioned decoder, Sinkhorn in none), so refitting Sinkhorn can only lower it \src{\file{docs/prereg/r4_a1_record.json}; SI-51; SI-57}.
\item The discriminator arm minimises scvi-tools' fool loss (cross-entropy to the uniform distribution over the other batches, over all cells) at one adversary step \src{SI-01; \file{scvi_adversarial_plan.py}:8--9, 292--298}; for two batches the target is the other batch alone, so this is the label-flipped non-saturating loss \src{\file{scvi_adversarial_plan.py}:297--298, with two batches}.
\item These results assume a fixed second distribution (the real data) and a moving generator (WGAN Thm.~1 and Example~1; Thms.~2.4--2.6), and Thm.~2.4 a discriminator close to the optimal one. Here every batch's latent moves and the discriminator takes one step per encoder step \src{SI-01; \file{scvi_adversarial_plan.py}:8--9; manifest}, so they do not transfer automatically; whether their conditions arise here is what the exploratory premise check R0 describes \src{SI-68}.
\item The vanishing-gradient premise is examined, not assumed: the premise check R0 (batch separability before integration; adversary gradient norms) is exploratory and gives no confirmed or refuted claim \src{SI-68; PAPER\_PLAN \S3; decision letter, Reviewer 1, major comment 2 (fairness of the comparison)}.
\end{itemize}

\paragraph{Composition shift and over-alignment.}
\begin{itemize}
\item When cell-type (label) proportions differ between two batches (domains), any one classifier applied to the shared latent space makes errors: if the label distributions are farther apart than the latent distributions, its error summed over the two batches is at least $\tfrac12\bigl(d_{JS}(\mathcal{D}^Y_S,\mathcal{D}^Y_T)-d_{JS}(\mathcal{D}^Z_S,\mathcal{D}^Z_T)\bigr)^2$, which grows as the latent distributions are aligned more tightly ($d_{JS}$: square root of the JS divergence; $\mathcal{D}^Y$, $\mathcal{D}^Z$: label and latent distributions of the two batches $S$, $T$) \citep[Thm.~4.3; binary labels]{zhao2019learning}.
\item Exact alignment with zero source error, under any distance that is zero only for identical distributions (a JS classifier, Wasserstein, maximum mean discrepancy (MMD)), leaves target error of at least the label-proportion gap \citep[Prop.~3.1; binary labels]{wu2019domain}.
\item Exact alignment under any such distance therefore costs accuracy under composition shift; whether the divergence arms over-align at the finite penalty strengths used here is a hypothesis, tested by the dose-response X3 \pending{X3 analysis}{specified during X1, before any X1 result is examined (SI-68)}.
\item An adversarial conditional variational autoencoder (VAE) and GLUE (a graph-linked embedding method) mixed unrelated cell types with unbalanced proportions, more so at higher correction strength \citep{hrovatin2025integrating}.
\item Dataset imbalance changed clustering, annotation, differential expression, mapping and trajectories for five methods \citep{maan2024characterizing}.
\end{itemize}

\paragraph{Loss versus tuning in GANs.}
\begin{itemize}
\item For image GANs, most losses reached similar scores with enough tuning and restarts; none consistently beat the non-saturating GAN \citep{lucic2018gans}.
\item The non-saturating loss was stable across datasets and hyperparameters; spectral normalisation consistently helped, the gradient penalty only sometimes \citep{kurach2019large}.
\item A critic--discriminator difference may thus reflect regularisation, update budget or tuning, not the divergence (image results) \pending{inference from image-GAN studies; examined for this setting by X6}{X6 analysis specified during X1 (SI-68)}.
\end{itemize}

\paragraph{Gap.}
\begin{itemize}
\item Closest controlled-backbone study: scIB-E, an extension of the single-cell integration benchmark scIB, with six batch losses (GAN, Hilbert--Schmidt independence criterion (HSIC), orthogonal projection, mutual information, gradient reversal, reverse cross-entropy) in one scVI framework; no W1, optimal-transport (OT) or explicit MMD term between batch distributions (its HSIC term measures dependence between latent code and batch label) \citep{yi2025benchmarking}.
\item No comparison of JS, W1 and MMD or OT objectives on one stock scVI backbone was found (searches of 2026-10-15 in OpenAlex, arXiv and Europe PMC from 2017, with citation tracking and web searches; scope and limits in \file{paper/evidence.md}): a search-based absence, not an exhaustive one \src{\file{paper/evidence.md}, Gap search} \pending{search repeated}{before submission}.
\item Reference anchoring and W1 critics already exist: a W1 critic for domain-invariant features \citep{shen2018wasserstein}, reference-anchored WGAN-GP integration \citep{wang2021imap} and MMD alignment to a reference or target batch \citep{amodio2019exploring,shaham2017removal}.
\end{itemize}

\paragraph{Contribution.}
\begin{itemize}
\item Scope: adversaries and divergences added to one stock scVI backbone at scvi-tools defaults; two decoders, conditioned (the decoder receives the batch label, as scVI is usually used) and unconditioned (no batch label reaches the model, so the adversary is the only batch correction), each with its own primary result \src{SI-09; SI-11; SI-16; SI-70}; equal epochs, minibatch and latent size within each experiment and task for every arm, plain scVI and sysVI, with the exceptions declared in Methods (Backbone) \src{SI-10; SI-04}; eight tasks of the scIB benchmark \src{SI-02; \file{docs/prepped_fingerprints_scib.json}}.
\item Limit of the primary contrast: P1 compares the reference critic at 5 critic steps per encoder step with the discriminator at 1 step \src{SI-01 (no 10-step discriminator); manifest}, so P1 alone does not separate the divergence from the extra critic steps; X6 varies the steps for the pooled critic only \src{manifest}.
\item A controlled decomposition of the critic--discriminator trade-off \pending{the decomposition as the paper's stated contribution (PAPER\_PLAN \S2; no ledger entry)}{your decision}. How much is:
\begin{enumerate}
\item operating point: arms compared at equal batch removal (P1, biology kept at equal batch removal) \src{SI-68};
\item composition-driven over-alignment: dose-response X3 (cell types depleted in one batch at rising doses, the same number of cells at every dose) \src{SI-25; SI-31--33; SI-41; manifest} \pending{X3 analysis}{specified during X1, before any X1 result is examined (SI-68)};
\item W1 geometry or latent scale: the W1 critics, MMD and Sinkhorn arms side by side in the single-strength tables and in the widths of the ranges of good strengths \src{SI-18; SI-68} \pending{attribution of a difference to W1 geometry rather than latent scale, and its analysis}{follow-up analyses locked during X1 (SI-68)}; per-dimension standardisation of the adversary input, compared in pilot A3 under rule R3 \src{SI-20; SI-27}; off in X1 for the same reason as the adversary input \src{SI-66; \file{docs/prereg/r4_a1_record.json}} \pending{formal re-run of rule R3}{after the Sinkhorn redo (SI-57)};
\item regularisation and update budget: X6 (the discriminator plain, spectrally normalised [\citealp{miyato2018spectral}] or with an R1 gradient penalty [\citealp{mescheder2018training}; not rule R1]; the pooled critic at 1, 5 and 10 critic steps, and spectrally normalised at 1 and 5) \src{SI-22; manifest} \pending{X6 analysis}{specified during X1 (SI-68)};
\item target design, reference versus symmetric: X7 \pending{X7 redesign (rated NO-GO as designed; final-review findings D-06 and I089)}{redesigned during X1, locked before any X1 result is examined (SI-68)}.
\end{enumerate}
\item Also asked: whether the critic--discriminator difference changes with the number of batches (X8) \pending{wording and design}{X8 redesign, during X1 (SI-68)}; whether the critic reaches matched batch removal through more local mixing \src{SI-76}; and how robust the results are to architecture, hyperparameters and KL warm-up (X12, X15) \src{SI-05; SI-44; PAPER\_PLAN \S4} \pending{analyses}{follow-up analyses locked during X1 (SI-68)}.
\item Comparison with stock scVI and standard tools (Harmony, Scanorama, sysVI, scANVI) at equal numbers of tries \src{SI-68; SI-77} \pending{outcome}{results of X1 and of the baseline experiment X13}. \pending{guidance for scVI users, including the case that no adversary beats stock scVI (PAPER\_PLAN \S2)}{results of X1 and X13}
\end{itemize}

% ---------------------------------------------------------------------------------------------------
% CLAIMS TO DROP OR SOFTEN (for the writer). Sources: PAPER_PLAN section 2; original_manuscript_claims_audit.csv
% (C-ids); this track's source checks (evidence.md). None may re-enter as settled.
%  - "JS fails on disjoint technology batches" / "vanishing gradients ... leading to failed training" (C01 partially
%    supported). R0 is exploratory (SI-68): it can describe separability and gradient norms, not confirm the premise.
%  - "First systematic comparison" / "no study has isolated the adversarial objective" (C02 untested). scIB-E compared six
%    batch losses inside scVI; WDGRL uses a W1 critic and iMAP WGAN-GP critics; iMAP and SAUCIE already use reference anchoring.
%  - "Critics over-correct" as a W1 property (C06 invalidated), unless shown at matched batch removal; Zhao 2019 (JS distance)
%    and Wu 2019 (exact alignment), both with binary labels, bound the error of aligned representations under composition shift;
%    sysVI reports that strong adversarial alignment mixed unrelated cell types in its methods.
%  - "Weak integration is protective" (C09 invalidated): that is the trade-off curve itself, unless shown at matched alignment.
%  - All topology language: "Topology Matters", "topological bottleneck", "radial projection", "topologically dense
%    reference", "topological collapse" (C12 untested, C14 contradicted; Reviewer 2, main concern 3: the work does not justify
%    topological mechanisms; PAPER_PLAN: nothing topological was measured).
%  - "W1 gives robust gradients regardless of initial separation" (C23): true of the distance (WGAN Thm 1), not of a
%    gradient-penalty critic, which enforces the Lipschitz constraint only softly (WGAN-GP Sec. 4).
%  - "Superior local mixing" (C04, C05), "scales poorly" (C10, C11), "broad lambda windows" (C16, C17), "~24% more
%    training time" (C18 contradicted), "ideal for severe disjoint batch effects" (C20 untested), barycenter as "the
%    theoretical solution" (C21), "primary limitation is over-correction" (C22).
%  - IMAAE as "Wasserstein-based" and as "motivating their decision to enforce a global Gaussian prior" (C24): the IMAAE
%    text never uses the word Wasserstein (its loss has WGAN-GP form) and the motivation link is the July authors' reading.
%  - "The unbalanced scenario is known to drive severe over-correction" (C25, uncited): cite sysVI and Maan 2024
%    for what they show, scoped to their methods; "balanced" was not composition-independent (C08 contradicted).
%  - July background attributions that the sources do not support (this track): batch effects from "the sequencing
%    technology" cited to Tung 2017 (Tung studies separately processed C1 replicates on one platform); the
%    statistical / geometric / deep-generative taxonomy cited to scIB (scIB groups methods by output: graph, embedding,
%    corrected matrix); "sample age" cited to Stegle 2015 (not checked; dropped; AD-AE names age as a confounder).
%  - PAPER_PLAN section 2 gap wording "trained on real counts": SI-15 fits the scIB files as distributed (no strict-count
%    main run), so "real counts" is dropped from the gap line.
%  - PAPER_PLAN section 3 "the premise becomes a tested hypothesis": SI-68 makes R0 exploratory.
%  - Pilot hint "stock scVI 0.726 vs best conditioned arm 0.646" (PAPER_PLAN section 2): provisional (one dataset, three
%    seeds, pre-stock schedule); not a result.
%  - August-revision sentence "the barycenter has not previously been used as a learnable ... target inside an
%    adversarial critic": a 'first' claim; not carried.
%  - July sentences "we implement a multi-headed critic formulation" / "We introduce a multi-class reference-based Wasserstein
%    critic" (C03 partially supported): here the reference critic is one of three critic formulations (reference, pooled,
%    barycenter); state the implemented generator objective exactly (claims audit C03).
% ---------------------------------------------------------------------------------------------------

\subsection{Research questions}
All six research questions are open: 19 of the 3,000 X1 rows are scored (local run paused on 2026-10-12) and the X1 analysis has not been run \src{NB-20261012-05; NB-20261014-21}. \pending{that no X1 result has been examined}{the user's confirmation at the diff review (ANALYSIS\_SPEC App.~A.3, item 2)}

\paragraph{R0 Premise (exploratory).}
\begin{itemize}
\item Question: before integration, are batches nearly separable in the plain-scVI latent, and does the discriminator's gradient vanish early in training while the critics' gradients persist? \src{PAPER\_PLAN \S3, R0 row; SI-68; ANALYSIS\_SPEC \S0}
\item Hypothesis: none. R0 is exploratory and makes no confirmed or refuted claim, so the premise that Jensen--Shannon (JS) gradients vanish is not stated as shown \src{SI-68; ANALYSIS\_SPEC App.~A.1, Q1 option A}.
\item Status: open.
\end{itemize}

\paragraph{R1 Integration versus conservation.}
\begin{itemize}
\item Question: at the same amount of batch removal, does the W1 (Wasserstein-1) critic keep a different amount of biology than the discriminator, and does it remove batch effects more locally? \src{PAPER\_PLAN \S3, R1 row; SI-76}
\item Hypothesis 1, two-sided: at the matched batch score b*, the reference critic and the discriminator (one adversary step per training step) differ in biology score C; tested once per decoder as P1, with its trade-off area \src{SI-01; SI-68; SI-70; SI-72; SI-74}.
\item Hypothesis 2, one signed question: at b*, the critic reaches the target through more local mixing, i.e.\ a higher mean of iLISI, kBET and graph connectivity than the discriminator \src{SI-76}.
\item Also tested: the collapse rate P2, critic minus discriminator \src{SI-75}. \pending{X3 composition dose-response hypothesis}{follow-up analyses locked during X1, before any X1 result is examined (SI-68)}
\item Status: open.
\end{itemize}

\paragraph{R2 Scalability with batch count.}
\begin{itemize}
\item Question: does the critic-minus-discriminator difference change with the number of batches? \src{topic: SI-68 (P3, batch-count scaling; ANALYSIS\_SPEC \S0)} \pending{exact wording (PAPER\_PLAN \S3, R2 row)}{X8 redesign, during X1 (SI-68)}
\item Hypothesis and endpoint P3: \pending{P3 wording and the X8 design (as designed, X8 cannot hold per-batch cell numbers and cell-type composition fixed across batch counts: D-07, I015, I077, I078)}{X8 redesign, locked during X1 before any X1 result is examined (SI-68)}
\item Status: open.
\end{itemize}

\paragraph{R3 Sensitivity to penalty strength and design factors.}
\begin{itemize}
\item Question: is any critic-versus-discriminator difference due to the divergence rather than to regularisation, update budget, latent scale or architecture, and how wide is each method's range of good penalty strengths? \src{PAPER\_PLAN \S3, R3 row; SI-68; ANALYSIS\_SPEC \S0}
\item Range of good strengths: measured on X1 with the selection score $S = 0.4B + 0.6C$ \src{SI-68; SI-77}; \pending{window rule (ANALYSIS\_SPEC \S10)}{diff review}.
\item Factor effects (X6, X12, X15, A3): \pending{hypotheses and analysis}{follow-up analyses locked during X1 (SI-68)}
\item Status: open.
\end{itemize}

\paragraph{R4 Reference selection.}
\begin{itemize}
\item Question: does the choice of reference batch matter more for the reference W1 critic than for reference-anchored JS or MMD with the same target? \src{topic: SI-68 (P4, reference choice; ANALYSIS\_SPEC \S0)} \pending{whether the question names cell-type coverage as the route (I089: X7 as designed supports an association only)}{X7 redesign, during X1 (SI-68)}
\item Hypothesis and endpoint P4: \pending{P4 wording; as designed, X7 supports at most an association with coverage (I089), and its variant arms are not at matched batch removal (D-06)}{X7 redesign, locked during X1 before any X1 result is examined (SI-68)}
\item Status: open.
\end{itemize}

\paragraph{R5 Placement against standard integrators.}
\begin{itemize}
\item Question: does any adversary improve on plain scVI, and where do the arms sit relative to Harmony, Scanorama, sysVI and scANVI when every method gets the same number of tries? \src{PAPER\_PLAN \S3, R5 row; SI-68; ANALYSIS\_SPEC \S0}
\item Locked analyses: tables of each method at one penalty strength chosen on the other datasets, and the comparison at equal numbers of tries, both ranked by $S = 0.4B + 0.6C$ \src{SI-68; SI-77}. \pending{rule for the single strength (leave-one-family-out) and best-of-k rule}{diff review (ANALYSIS\_SPEC \S8, \S9)}
\item Hypothesis: \pending{whether R5 states a hypothesis (PAPER\_PLAN poses it as a question)}{your decision}
\item Status: open.
\end{itemize}

% Taxonomy carried over from the July submission (statistical, geometric, deep generative, adversarial):
%   Topology_Matters.zip oup-authoring-template.tex:137,139 = 08f1e3c manuscript_tm/oup-authoring-template/
%   oup-authoring-template.tex:135,137; additions named in PAPER_PLAN section 3, item 2. One line per work.

\section{Related work}\label{sec:related}

\paragraph{Statistical and geometric integration.}
\begin{itemize}
\item limma: gene-wise linear models that can include batch terms; \texttt{removeBatchEffect} \citep{ritchie2015limma}.
\item ComBat: empirical Bayes batch adjustment, robust for small batches \citep{johnson2007adjusting}.
\item MNN: correction from mutual nearest neighbours between batches \citep{haghverdi2018batch}.
\item Seurat: anchors between datasets, across technologies and modalities \citep{stuart2019comprehensive}.
\item Harmony: shared embedding where cells group by cell type, not dataset \citep{korsunsky2019fast}.
\item Scanorama: finds and merges the cell types shared by pairs of datasets \citep{hie2019efficient}.
\end{itemize}

\paragraph{Benchmarks.}
\begin{itemize}
\item 14 methods, five scenarios, including non-identical cell types across batches \citep{tran2020benchmark}.
\item scIB: 68 method and preprocessing combinations, 13 tasks, 14 metrics of batch removal and biological conservation \citep{luecken2022benchmarking}; eight of its 13 tasks are used here \src{SI-02; \file{docs/prepped_fingerprints_scib.json}}.
\item Open Problems: benchmarks with different datasets and metrics disagree; standardised tasks \citep{luecken2025defining}.
\item scIB-E (extended scIB): six losses on the batch label in one scVI framework (GAN, HSIC, orthogonal projection, mutual-information minimisation, gradient reversal, reverse cross-entropy), then label-aware levels in scANVI; loss weights tuned with Ray Tune (automated hyperparameter search) on the immune task \citep{yi2025benchmarking}.
\end{itemize}

\paragraph{Deep generative models.}
\begin{itemize}
\item scVI: zero-inflated negative binomial model conditioned on batch, with a low-dimensional Gaussian latent \citep{lopez2018deep}; implemented in scvi-tools \citep{gayoso2022python}.
\item scANVI (single-cell annotation using variational inference): semi-supervised scVI that uses cell-type labels \citep{xu2021probabilistic}.
\item sysVI: conditional VAE with a VampPrior (a multimodal variational mixture-of-posteriors prior) and latent cycle-consistency, for substantial batch effects \citep{hrovatin2025integrating}.
\item trVAE (transfer VAE): conditional out-of-distribution generation for unpaired data, for example perturbation response; MMD on the decoder layer that follows the bottleneck matches the conditions \citep{lotfollahi2020conditional}.
\end{itemize}

\paragraph{Adversarial (JS-type) integration.}
\begin{itemize}
\item DANN (domain-adversarial neural network): domain-indiscriminate features via gradient reversal against a domain classifier \citep{ganin2016domain}.
\item AD-AE (adversarial deconfounding autoencoder): autoencoder plus an adversary that predicts the confounder \citep{dincer2020adversarial}.
\item scDREAMER (single-cell deep generative integration model with adversarial classifier): adversarial VAE (Bhattacharyya loss) plus an adversarial cross-entropy batch classifier \citep{shree2023scdreamer}.
\item scvi-tools' adversarial batch classifier: fool loss weighted by one minus the KL weight (arm scvi\_adv) \src{scvi-tools 1.4.2: scvi/train/\_trainingplans.py:535, 674--695; \file{scvi_adversarial_plan.py}:175--177}.
\end{itemize}

\paragraph{Wasserstein critics.}
\begin{itemize}
\item WDGRL (Wasserstein distance guided representation learning): a gradient-penalised domain critic estimates the empirical W1 between source and target features, which the feature extractor minimises \citep{shen2018wasserstein}.
\item iMAP (integration of multiple single-cell datasets by adversarial paired transfer networks): adversarial losses optimised with WGAN-GP; one dataset is the uncorrected anchor \citep{wang2021imap}.
\item ResPAN: WGAN with random-walk mutual-nearest-neighbour pairing and skip-connected autoencoders \citep{wang2022respan}.
\item IMAAE (integrating multiple single-cell datasets via an adversarial autoencoder): adversarial autoencoder matching the latent to $\mathcal{N}(0,I)$ with a WGAN-GP-form loss (never named Wasserstein), an anchor batch and known cell labels \citep{wang2023integrating}.
\item In the IMAAE paper, a GAN-only model tried at the design stage under-corrected with three or more batches (one dataset, judged from UMAP plots); the authors then used the GAN to constrain the latent space of an autoencoder \citep{wang2023integrating}.
\end{itemize}

\paragraph{Kernel and optimal-transport alignment.}
\begin{itemize}
\item MMD: largest difference in expectations over the unit ball of a reproducing-kernel Hilbert space \citep{gretton2012kernel}.
\item MMD-ResNet: residual network minimising the MMD between one source batch and one target batch \citep{shaham2017removal}.
\item SAUCIE (sparse autoencoder for unsupervised clustering, imputation and embedding): MMD penalty on an internal layer aligning batches to a reference batch \citep{amodio2019exploring}.
\item BERMUDA (batch effect removal using deep autoencoders): MMD between similar clusters across batches; tolerates different compositions if two batches share a cell type \citep{wang2019bermuda}.
\item uniPort: coupled VAE with minibatch unbalanced OT \citep{cao2022unified}.
\item Entropy-regularised OT \citep{cuturi2013sinkhorn}; Sinkhorn divergences interpolate between OT and MMD \citep{genevay2018learning,feydy2019interpolating}; Wasserstein barycenters \citep{cuturi2014fast}; POT (Python Optimal Transport) library \citep{flamary2021pot}.
\item \pending{further single-cell methods from the review's coverage search, read at abstract level only: iDLC (optimal-transport-regularised adversarial alignment, 2026 preprint), scStarCorrect (adversarial alignment to one fixed reference batch), Beaconet (adversarial, no reference), Portal (adversarial domain translation), AIF, ABC and IMGG (adversarial), and the batch-removal versus biology trade-off curve on scVI of Li et al. 2022}{your decision on coverage; each line then checked in the full text}
\end{itemize}

\section{Methods}
Reproduction detail is in the Supplementary material [S-...].

\subsection{Backbone}\label{m:backbone}
[S-M1, S-E2, S-E3]
\begin{itemize}
\item scvi-tools SCVI (\citealp{lopez2018deep}; \citealp{gayoso2022python}) at its defaults for every arm, for plain scVI and for the scVI pretraining of scANVI: 10 latent dimensions, one hidden layer of 128 units, zero-inflated negative binomial (ZINB) likelihood, 90\% training split, minibatch 128 \src{SI-10; \file{build_paper_manifest.py}:62; \file{fit_paper_config.py}:393--394; manifest}; sysVI uses the same latent size, hidden width, layers, training split, epochs and minibatch with its own Gaussian likelihood on the normalised values \src{\file{fit_paper_config.py}:404--413; SI-63}.
\item Epochs by the scvi-tools heuristic $\min(400, \mathrm{round}(20000 \times 400/n))$ on the $n$ cells fitted, no early stopping; within an experiment and task every arm, plain scVI and sysVI share epochs, minibatch and latent size; declared exceptions: scANVI's extra epochs, X12's factorial levels, and the subsample-sized epochs of X3 and X8 \src{SI-04; \file{build_paper_manifest.py}:114; \file{scvi_adversarial_plan.py}:563; \file{fit_paper_config.py}:401; PREREG \S9, 2026-10-04 (CR-12); manifest}.
\item Scope: scVI as people use it \src{SI-09}.
\item Plain scVI (arm none, $\lambda = 0$) is the stock training plan; a committed test asserts bit-identity with stock SCVI \src{\file{scvi_adversarial_plan.py}:4; \file{tests/scvi/test_adversarial_plan_v2.py}:48}.
\item Input: the counts layer on the file's highly variable genes; output: the posterior mean \src{\file{scvi_adversarial_plan.py}:528; \file{fit_paper_config.py}:321; \file{scvi_adversarial_plan.py}:568}.
\end{itemize}

\subsection{Decoders}
[S-M1]
\begin{itemize}
\item Conditioned: the decoder receives the batch label. Unconditioned: no batch label reaches encoder or decoder, only the adversary \src{\file{scvi_adversarial_plan.py}:534, :537; \file{tests/scvi/test_adversarial_plan_v2.py}:118}.
\item Both decoders in X1, 5 seeds each \src{SI-11; SI-16; manifest}.
\end{itemize}

\subsection{Adversarial arms of X1}
[S-M1--S-M7, S-E2]
\begin{itemize}
\item Each arm except scvi\_adv adds $\lambda \times$ (alignment term) to the loss, $\lambda$ fixed throughout \src{\file{scvi_adversarial_plan.py}:436, :177; \file{fit_paper_config.py}:340; manifest}.
\item Discriminator (JS-type; the minimax objective of \citealp{goodfellow2014generative}, stated there for two distributions): classifies each cell's batch; the encoder minimises scvi-tools' bounded fool loss; 1 discriminator step per training step \src{SI-01; \file{scvi_adversarial_plan.py}:206, :292; manifest}.
\item Reference critic (W1; \citealp{arjovsky2017wasserstein}): one head per batch estimates W1 to a reference batch, whose cells also move; 5 critic steps \src{\file{wcd/critic.py}:72; \file{scvi_adversarial_plan.py}:421; manifest}.
\item Pooled critic (W1): head $k$ targets all other batches \src{\file{wcd/critic.py}:77}.
\item Barycenter critic (W1): head $k$ targets a Wasserstein-2 barycenter (average distribution; \citealp{cuturi2014fast}) of the minibatch's batches, recomputed per step, cold start, at most 10 iterations \src{SI-14 (10 iterations); \file{wcd/barycenter.py}:38, :72--73 (cap 10, stop threshold $10^{-7}$); manifest}.
\item MMD (no network): squared maximum mean discrepancy (\citealp{gretton2012kernel}) of each batch to the others, fixed multi-scale Gaussian kernel \src{\file{wcd/alignment.py}:49, :109}.
\item Sinkhorn (no network): debiased Sinkhorn divergence of each batch to the others (\citealp{cuturi2013sinkhorn}; \citealp{genevay2018learning}; \citealp{feydy2019interpolating}), Euclidean cost, entropic regularisation $\varepsilon$ = blur = $0.05\sqrt{2d}$ \src{\file{wcd/alignment.py}:12--13, :50, :249}; not a W1 estimate at this blur \src{\file{wcd/alignment.py}:15}; \pending{size of the gap to W1}{recomputed on refitted latents}.
\item scvi\_adv: scvi-tools' own classifier adversary (32 hidden units, weight $1 -$ KL weight), conditioned decoder only, outside rules R1--R4 \src{\file{scvi_adversarial_plan.py}:5, :177; PREREG \S3; manifest}.
\item Discriminator and critics: same head architecture (latent, 128, 128, one output per batch; ReLU) \src{\file{wcd/adversarial.py}:50, :76}; critics add a per-head gradient penalty (\citealp{gulrajani2017improved}) of weight 10 \src{\file{wcd/critic.py}:158}.
\item Reference batch: highest cell-type Shannon entropy, ties by size \src{\file{wcd/data.py}:32; \file{fit_paper_config.py}:324; manifest}.
\end{itemize}

\subsection{Adversary input and standardisation}
[S-M8, S-M11]
\begin{itemize}
\item Defaults: posterior-mean input, no per-dimension standardisation; rules R2 and R3 switch either only if their criteria hold, and a sample input that misses b* where the mean reaches it is refused \src{SI-19; SI-20; PREREG \S3}.
\item One pooled decision over both decoders (A2: 12 cells, 8 evaluable needed; A3: 20, 14) \src{SI-27}; low-signal cells count as not evaluable \src{SI-66}; edge window points are fitted on A1 seeds \src{SI-45}.
\item The decision covers all six arms; if both switch, the untested combination is disclosed \src{PREREG \S3}.
\item X1 keeps the defaults (posterior mean; standardisation off): the frozen rule-R4 record on A1 flags immune with the conditioned decoder as low-signal ($b_{\mathrm{common}} - b_0 = 0.0188$; sim1 with the conditioned decoder is also flagged, 0.0069, but is not an A2 or A3 task), so rules R2 and R3 cannot switch while that flag holds \src{SI-66; \file{docs/prereg/r4_a1_record.json}; NB-20261012-01}; Sinkhorn sets $b_{\mathrm{common}}$ in no block, so refitting it alone can only lower $b_{\mathrm{common}}$ \src{\file{docs/prereg/r4_a1_record.json} targets}. \pending{formal re-run of rules R2 and R3}{after the Sinkhorn redo (SI-57)}
\end{itemize}

\subsection{Training}
[S-E2, S-E3, S-M6]
\begin{itemize}
\item Generator: scvi-tools' Adam (learning rate $10^{-3}$, $\epsilon$ 0.01, weight decay $10^{-6}$) \src{\file{scvi_adversarial_plan.py}:245}. Discriminator: the same Adam settings, 1 step. Critics: Adam $10^{-4}$, betas (0, 0.9), 5 steps \src{\file{scvi_adversarial_plan.py}:258, :260; manifest; PREREG \S9, 2026-10-04 (CR-12)}.
\item Adversary steps reuse the minibatch's detached latent (WGAN-GP draws fresh data) \src{\file{scvi_adversarial_plan.py}:392--394; PREREG \S9, 2026-10-04 (CR-12)}.
\item KL warm-up: scvi-tools' 400-epoch default, incomplete on tasks with fewer epochs [S-E3], while $\lambda$ acts from the first step \src{\file{fit_paper_config.py}:123; manifest; PREREG \S9, 2026-10-04 (X15)}.
\item Failed fit: non-finite training loss (diverged) or non-finite latent (nonfinite\_latent); infrastructure errors are rerun \src{PREREG \S1}.
\item Sinkhorn corrected before X1 (gradient counted once; solve converged at the final smoothing step) \src{SI-51}; its 180 A1 fits are refitted, then rules R1 and R4 on A1 decide which extension and A3 rows to refit (18 and 36 before the fix, 234 fits in all) \src{SI-51; SI-57}.
\item Sinkhorn solve in float64, at most 500,000 final-step iterations with an error at the cap, immediate stop on a NaN marginal error \src{SI-61; SI-62}. \pending{tolerance ($10^{-4}$ L1 marginal error, \file{wcd/alignment.py}:56) and the outcome of a fit that keeps reaching the cap}{diff review, tag prereg-tier12-v3}
\end{itemize}

\subsection{Data}
[S-D1, S-D2]
\begin{itemize}
\item Eight of the scIB benchmark's integration tasks (\citealp{luecken2022benchmarking}), no new datasets: six real (pancreas, lung, immune, immune\_hum\_mou, and the gene-activity versions of atac\_small and atac\_large) and two simulations (sim1, sim2) \src{SI-02; \file{build_paper_manifest.py}:54; \file{docs/prepped_fingerprints_scib.json}}; scIB's batch and label keys \src{SI-03}.
\item Fitted as distributed, with the counts layer exactly as scIB fed scVI \src{SI-15}; 18 batches in 4 tasks hold non-integer values \src{\file{docs/prepped_fingerprints_scib.json}}.
\item Genes: 2,000 highly variable genes per RNA task (scib.pp.hvg\_batch); all features of the ATAC gene-activity matrices \src{\file{docs/prepped_fingerprints_scib.json}}.
\item Families: pancreas; lung; immune (immune, immune\_hum\_mou); ATAC (atac\_small, atac\_large) \src{SI-73}.
\item immune\_hum\_mou contains all 33,506 immune cells \src{\file{D/prepped_scib} obs names}.
\item Every one of the 11,270 atac\_small cells is also in atac\_large (same batch, barcode and cell-type label; 3,618 of them, the 10x Genomics batch, under identical names); count values differ between the two files (0 of 200 sampled matched pairs identical) \src{\file{D/prepped_scib} obs and counts layers}.
\item Simulations are not pooled with real tasks \src{PREREG \S5}.
\end{itemize}

\subsection{Evaluation}
[S-E2, S-E6]
\begin{itemize}
\item scib.metrics.metrics (scib 1.1.7) on the posterior-mean embedding and its neighbour graph \src{SI-08; \file{score_scib_native.py}:207, :214}.
\item B: mean of PCR comparison, batch ASW, graph iLISI, graph connectivity, kBET (\citealp{buttner2019test}) \src{PREREG \S1; \file{score_scib_native.py}:51}. C: mean of NMI, ARI, cell-type ASW, isolated-label F1 and silhouette, graph cLISI, plus cell-cycle conservation (pancreas, lung, immune, immune\_hum\_mou) and trajectory (immune, immune\_hum\_mou, the tasks with pseudotime) \src{PREREG \S1; SI-28; \file{score_scib_native.py}:52, :212--213}.
\item Raw values on $[0, 1]$; no rule uses min--max scaling; required metrics must be finite \src{PREREG \S1}.
\item kBET: R seed 0 set before every call; rule scripts refuse mixed kBET seeds \src{PREREG \S1; \file{score_scib_native.py}:54, :80}.
\item NMI of 0 at every resolution keeps scib's error, decided case by case \src{SI-65}; latents without cell names are refused \src{SI-64}.
\item Subsets (X3, X8) use the all-feature PCA reference; score rows record scorer commit, CPU and versions, and mixed provenance is refused \src{PREREG \S9, 2026-10-04 (CR-01, CR-05, CR-08, CR-14)}.
\item Sensitivity: X1 comparisons on immune and immune\_hum\_mou also with C without trajectory \src{SI-38}.
\end{itemize}

\subsection{Pre-registered design}
[S-M10--S-M12, S-E1, S-E4, S-E6]
\begin{itemize}
\item Tier 1 and Tier 2 experiments (Tier 2 adds X3, X6, X12) \src{SI-05; PAPER\_PLAN \S4}; all strength grids re-derived from new pilot fits, none selected from earlier sweeps \src{SI-13}.
\item Stage order: A1, rule R1, rule R4 on A1, A2/A3, rules R2/R3, X1, rule R4 on X1, follow-ups; rule-free fits (X1 none and scvi\_adv, X13) may run earlier \src{PREREG \S0}.
\item Seeds: A1--A3 100--102, X1 and X13 0--4, follow-ups 10--12, so no follow-up reuses an X1 fit that chose its strength \src{SI-26; PREREG \S4}.
\item Rule R1: one six-point grid per family and decoder (8 grids), from no effect to over-correction or collapse \src{SI-16; SI-18; PREREG \S2}; constants: smallest effect 0.01, noise multiple 3, collapse 0.10 \src{PREREG \S1}.
\item Rule R4: $b^* = b_0 + 0.5\,(b_{\mathrm{common}} - b_0)$ per task and decoder, unscaled; matched strength = failure-free grid point with seed-mean B closest to b* \src{SI-21; PREREG \S4}.
\item Grid edges: extensions only for matching; P1 and the frontier use the six registered strengths; edge optima disclosed \src{SI-55; PREREG \S4}.
\item Curves of C against B start at $\lambda = 0$; low-signal blocks ($b_{\mathrm{common}} - b_0 < 0.02$) are not evaluable for P1 or rules R2/R3 and are reported separately \src{SI-56; SI-66; PREREG \S4}.
\item Freeze: rule files and manifest builder tagged before first A1 fit; each rule's record (input hashes, commit) committed before its stage \src{PREREG \S6}. Manifest of record: 6,930 rows, SHA-256 6efbe422c8f09f98\ldots{} \src{PREREG \S9, 2026-10-04; recomputed}.
\item Fits and scores at c88cce3 and at the new tag form one set except for Sinkhorn, whose rows must carry the new tag; no rescoring \src{SI-54; SI-60; \file{docs/EQUIVALENT_COMMITS.json}}.
\item \pending{freeze of version 3 (amendments of 2026-10-13/14, ANALYSIS\_SPEC, \file{analysis_x1.py})}{diff review, then tag prereg-tier12-v3 (SI-52, SI-60)}
\end{itemize}

\subsection{Statistical analysis of X1}
[S-M13--S-M18]
\begin{itemize}
\item Locked now: P1 with its trade-off area, P2, local versus global, single-strength tables, equal-tries comparison with standard tools, window width; R0 and X10 exploratory \src{SI-68}.
\item P1 contrast: reference critic minus discriminator at 1 step \src{SI-01; PREREG \S5; manifest}.
\item Average the 5 runs per strength, read C at the first crossing of b*; leave-one-seed-out uncertainty \src{SI-69}.
\item One primary result per decoder, never averaged \src{SI-70}.
\item Failed runs left out of P1 and the area, counted in P2 \src{SI-71}.
\item Area corner: plain scVI's B and its C minus 0.10; divided by the task's batch range to pool \src{SI-72}.
\item Pooling: overall + family + task effects, prior scales 0.1 (0.05, 0.2 as sensitivity), exact calculation \src{SI-73}.
\item Margin 0.01; direction, practical equivalence or inconclusive at posterior probability 0.95, the region-of-practical-equivalence rule of \citealp{corani2017statistical} \src{SI-74}.
\item P2: collapse = failed, or C at least 0.10 below the same seed's plain scVI; Clopper--Pearson intervals (\citealp{clopper1934use}); pooled critic-minus-discriminator difference (prior 0.25, margin 5 points) \src{SI-75}. \pending{whether SI-75 stands after the pilot recount (8 of 54, not 8 of 45)}{diff review}
\item Local versus global: mean of iLISI, kBET, graph connectivity at the P1 crossing; sign test plus the P1 model, margin 0.01 \src{SI-76}.
\item Selection score $S = 0.4B + 0.6C$, unscaled \src{SI-77}.
\item Frequentist companion: exact two-sided Wilcoxon signed-rank test on the per-task differences, with Holm's correction over the four primary endpoints P1--P4; with 6 real tasks its smallest p, $2/2^6 = 0.031$, is above Holm's first threshold $0.05/4 = 0.0125$, so it cannot reject at the Holm-adjusted level \src{PREREG \S5; \file{docs/prereg/r1_record.json} power}. \pending{Holm over six primary results (first threshold 0.0083) and reporting the companion as descriptive}{diff review (ANALYSIS\_SPEC \S4, \S5)}
\item \pending{agent proposals of ANALYSIS\_SPEC \S2.3, \S5, \S8, \S9, \S10, \S11 (unnamed sensitivities), \S15}{diff review; approvals become ledger entries}
\item \pending{definition of ``examined''; redone R1 record name; scaling sensitivities}{diff review (ANALYSIS\_SPEC App.~A.3)}
\item Power (PREREG \S5 rule; frozen record): the design is flagged UNDERPOWERED; minimal detectable effect at 80\% power 0.131 for D = 6 real tasks (0.182 for D = 4 families) [S-M15] \src{PREREG \S5; \file{docs/prereg/r1_record.json} power}. \pending{recomputed estimate with the P1 reading and evaluable blocks only}{rule R4 on A1 re-run after the Sinkhorn redo (SI-57)}
\item \pending{one analysis run after all 3,000 X1 and 408 X13 rows are scored or failed}{tag prereg-tier12-v3}
\end{itemize}

\subsection{Follow-up experiments}
[S-E1, S-M9]
\begin{itemize}
\item Seeds 10--12, at the strengths matched by rule R4 on X1; variant arms inherit their base arm's strength \src{PREREG \S4; SI-26}. \pending{matched strengths}{rule R4 on the scored X1}
\item X3: one non-reference batch depleted of two cell types (sim2: one) at doses 0--100\%, equal cell count per dose, reference fixed at the dose-0 choice, importance-weighted control at 50, 80, 95\% \src{SI-25; SI-31--33; SI-41; manifest}. \pending{draw and weights as implemented (SI-42, lead-agent decision)}{the user's confirmation}
\item X6: discriminator plain, spectrally normalised (\citealp{miyato2018spectral}) or with an R1 gradient penalty (\citealp[Eq.~9]{mescheder2018training}; $\gamma = 10$ is this study's setting); pooled critic at 1, 5 and 10 steps, and spectrally normalised at 1 and 5 steps \src{SI-22; manifest}.
\item X12: 8-run fractional factorial, resolution III, over latent size, adversary width and learning rate, minibatch and linear-or-default decoder \src{\file{build_paper_manifest.py}:259, :269; I028}.
\item X15: stock against completing KL warm-up, immune and atac\_large, 84 fits \src{SI-44; manifest}.
\item \pending{analyses of X3, X6, X12, X15}{locked during X1, before any X1 result is examined (SI-68)}
\item X7: \pending{redesign (arms not at matched batch removal, D-06; coverage confounded with reference size and entropy, I089)}{during X1, before any X1 result is examined (SI-68)}; its reference-anchored JS loss is fixed \src{SI-23}.
\item X8: \pending{redesign (per-batch cell numbers and cell-type composition change with batch count; D-07, I015, I077, I078)}{during X1, before any X1 result is examined (SI-68)}; its critic-only stratified sampler is fixed \src{SI-24}.
\end{itemize}

\subsection{Baselines (X13)}
[S-E1, S-E2]
\begin{itemize}
\item Plain scVI = X1 arm none \src{manifest}; scANVI adds $\min(10, \max(2, \mathrm{round}(E/3)))$ epochs after its scVI pretraining \src{SI-04; \file{fit_paper_config.py}:401; PREREG \S9, 2026-10-04 (CR-12); manifest}.
\item Harmony (\citealp{korsunsky2019fast}): $\theta \in \{0, 0.5, 1, 2, 4, 8\}$ at 10 PCs, once each, fresh process, pinned threads; 50 PCs at $\theta$ 2 as sensitivity \src{SI-34}.
\item Scanorama (\citealp{hie2019efficient}): knn $\in \{2, 5, 10, 20, 40, 80\}$ at 10 dimensions, once each; 100 dimensions at knn 20 as sensitivity \src{SI-46}.
\item PCA: once at 10 (and 50) components, the uncorrected anchor \src{SI-36}.
\item scANVI (\citealp{xu2021probabilistic}): scIB's protocol, scVI then scANVI with all labels; seeds 0--4 \src{\file{fit_paper_config.py}:399, :401; manifest}.
\item sysVI (\citealp{hrovatin2025integrating}): scvi-tools' SysVI, cycle weight $\in \{1, 2, 5, 10, 20, 50\}$, seeds 0--4, no KL warm-up \src{\file{fit_paper_config.py}:404, :413; manifest; PREREG \S9, 2026-10-04 (CR-12)}.
\item Harmony, Scanorama and PCA use the file's normalised values on its highly variable genes \src{\file{run_cpu_baselines.py}:230}.
\item \pending{placement rule (best-of-k at equal k; scANVI in a label-aware stratum)}{diff review (ANALYSIS\_SPEC \S9)}
\end{itemize}

\subsection{Compute and reproducibility}
[S-E5, S-E6, S-D2]
\begin{itemize}
\item Each task runs entirely on one machine \src{SI-17}.
\item atac\_small, immune, sim1: local RTX 3080 \src{SI-30}; the other five: JHPCE L40S GPUs \src{SI-39}.
\item Scoring is local, from one worktree pinned at the tag per stage \src{SI-47; PREREG \S9, 2026-10-04 (b)}.
\item Software [S-E5]: scvi-tools 1.4.2, PyTorch 2.13.0 (fits); scib 1.1.7, R 4.5.2, kBET 0.99.6 (scoring) \src{\file{cluster/envspec/local_scvi-}\file{api_pip_freeze.txt}; \file{docs/jhpce/gate_data/versions_*_local.json}}.
\item Each fit sets scvi-tools' seed to its row seed \src{\file{scvi_adversarial_plan.py}:526}; each latent stores its row, resolved settings and provenance \src{\file{fit_paper_config.py}:4--8}.
\item Identity checks at c3fd92d (CPU, 1 epoch, 1 thread, one fit per code path): 59 of 59 non-Sinkhorn fit code paths and 7 of 7 rescored latents (all 32 compared columns) equal c88cce3 \src{\file{docs/EQUIVALENT_COMMITS.json}; NB-20261013-39, -40}.
\item From c3fd92d to 08f1e3c no file under \file{src/} and none of \file{fit_paper_config.py}, \file{scvi_adversarial_plan.py}, \file{score_scib_native.py} or their imports changed, which is the validity condition of \file{docs/EQUIVALENT_COMMITS.json}; the stage runner \file{run_stage.py}, which is in the fit import closure, changed its start-up provenance check (CF-02, 8351999), and the proof was taken to carry over without a re-run \src{git log c3fd92d..08f1e3c; NB-20261014-29 (agent judgment)}. \pending{acceptance of the carry-over; tag commit, with the checks repeated if fit or score code changes}{diff review, tag prereg-tier12-v3 (SI-52)}
\item Data provenance: scIB source files and MD5 checksums [S-D2] \src{\file{docs/prepped_fingerprints_scib.json}}.
\end{itemize}

\section{Results}
PAPER\_PLAN \S3 lists nine figures; that list is a proposal.

\subsection{R0 Premise (exploratory)}
\begin{itemize}
\item Endpoint and rule: none; exploratory, no confirmed or refuted claim \src{SI-68}.
\item Data: \pending{X1 training histories and $\lambda = 0$ latents (ANALYSIS\_SPEC \S0)}{your decision}
\item Figure: \pending{premise figure}{your decision}
\item Result: \pending{R0 result}{X1 scored and the single analysis run}
\end{itemize}

\subsection{R1 Integration versus conservation}
\begin{itemize}
\item Endpoints: P1 per decoder with its trade-off area [S-M13, S-M14], P2 [S-M16], local versus global [S-M17] \src{SI-68; SI-70; SI-72; SI-75; SI-76}.
\item Rule: declaration on the pooled posterior over the 6 real tasks, margin 0.01 (P2: 5 points); low-signal blocks reported separately; simulations per task \src{SI-73; SI-74; SI-75; SI-56; PREREG \S5}.
\item Experiments: X1 \src{SI-68}; X10, exploratory \src{SI-68}; X3 \pending{composition dose-response analysis}{locked during X1 (SI-68)}
\item Figure: \pending{frontier and P1 figures (PAPER\_PLAN \S3 proposes per-task frontiers and a forest plot)}{your decision}
\item Result: \pending{P1, area, P2 and local-versus-global results per decoder}{X1 scored and the single analysis run}
\end{itemize}

\subsection{R2 Scalability with batch count}
\begin{itemize}
\item Endpoint and rule: \pending{P3 and its decision rule}{X8 redesign (SI-68)}
\item Experiments: X8 \pending{redesigned X8}{redesign during X1}
\item Figure: \pending{scaling figure}{your decision}
\item Result: \pending{P3 result}{X8 fitted and scored, follow-up analysis run}
\end{itemize}

\subsection{R3 Sensitivity to penalty strength and design factors}
\begin{itemize}
\item Endpoint (X1 part): width of each method's range of good strengths, ranked by S [S-M18] \src{SI-68; SI-77}; \pending{window rule}{diff review (ANALYSIS\_SPEC \S10)}
\item Factor effects: \pending{endpoints and rules for X6, X12, X15 and A3}{locked during X1 (SI-68)}
\item Figure: \pending{strength-curve and factor figures}{your decision}
\item Result: \pending{window widths; factor effects}{X1 scored and the single analysis run; follow-up analysis run}
\end{itemize}

\subsection{R4 Reference selection}
\begin{itemize}
\item Endpoint and rule: \pending{P4 and its decision rule}{X7 redesign (SI-68)}
\item Experiments: X7 \pending{redesigned X7}{redesign during X1}
\item Figure: \pending{reference-coverage figure}{your decision}
\item Result: \pending{P4 result}{X7 fitted and scored, follow-up analysis run}
\end{itemize}

\subsection{R5 Placement against standard integrators}
\begin{itemize}
\item Endpoints: single-strength tables and the comparison at equal numbers of tries, ranked by S \src{SI-68; SI-77}; PCA is the uncorrected anchor \src{SI-36}; \pending{transfer and best-of-k rules}{diff review (ANALYSIS\_SPEC \S8, \S9)}
\item Experiments: X13 and X1 \src{SI-68}.
\item Figure: \pending{placement figure}{your decision}
\item Result: \pending{placement and single-strength results}{X1 and X13 scored and the single analysis run}
\end{itemize}

\section{Discussion}
\subsection{Interpretation}
\begin{itemize}
\item \pending{what P1, its trade-off area, P2 and the local-versus-global result mean, per decoder}{X1 scored and the single analysis run}
\item \pending{guidance for scVI users, including the case that no adversary improves on plain scVI (PAPER\_PLAN \S3)}{X1 and X13 analysed}
\item \pending{mechanisms: composition shift, regularisation and update budget, reference design, batch count}{follow-up analyses locked and run (SI-68)}
\end{itemize}

\subsection{Limitations}
\begin{itemize}
\item One categorical batch variable per task, scIB's batch key \src{SI-03; \file{docs/prepped_fingerprints_scib.json}}.
\item The ATAC tasks are gene-activity matrices, not peak counts \src{\file{docs/prepped_fingerprints_scib.json}}.
\item Strength ranges come from 3 of the 8 tasks (atac\_small, immune, sim1), fitted on seeds disjoint from X1 \src{PREREG, Disclosure; PREREG \S8}.
\item Eighteen batches with non-integer values (pancreas, lung, both immune tasks) are fitted as distributed under a count likelihood \src{SI-15; \file{docs/prepped_fingerprints_scib.json}}.
\item Simulations are a separate family, not pooled with real tasks \src{PREREG \S5}.
\item Six real tasks in four families, two of them nested: immune\_hum\_mou contains all immune cells, and atac\_large contains all atac\_small cells, with different count values \src{SI-73; \file{D/prepped_scib} obs and counts layers}.
\item The design is underpowered (minimal detectable effect 0.131 for 6 real tasks at 80\% power), and the Wilcoxon companion cannot reject at the Holm-adjusted level (smallest p 0.031 $>$ 0.0125) \src{PREREG \S5; \file{docs/prereg/r1_record.json} power}. \pending{recomputed power and minimal detectable effect}{rule R4 on A1 re-run (SI-57)}
\item P1 compares the reference critic at 5 critic steps with the discriminator at 1; an update-budget control exists only for the pooled critic (X6) \src{SI-01; manifest}.
\item The KL warm-up ends at weight 0.20--0.61 on the four tasks with fewer than 400 epochs, while $\lambda$ acts from the first step \src{\file{fit_paper_config.py}:123; manifest}; X15 tests this \src{SI-44}.
\item Adversary steps reuse one minibatch \src{\file{scvi_adversarial_plan.py}:392--394}.
\item Scoring is transductive: the scored embedding covers all cells, including the 90\% used for training \src{\file{scvi_adversarial_plan.py}:568; manifest}.
\item Rule R1's noise estimate depends on the trajectory metric; R1 is kept as registered and also reported without trajectory \src{SI-53}.
\item The Sinkhorn arm is not a W1 estimate at its blur \src{\file{wcd/alignment.py}:15}.
\end{itemize}


\renewcommand{\bibsection}{\section*{References}\addcontentsline{toc}{section}{References}}
\renewcommand{\bibfont}{\raggedright}  % long proceedings URLs otherwise leave empty justified lines
\bibliographystyle{plainnat}
\bibliography{refs}


\clearpage
\section*{Supplementary material}
Equations are read from the code at commit 08f1e3c; file:line as in the reading guide.

\subsection*{S-M Methods detail}
\paragraph{S-M1 Generator objective, backbone loss and KL weight.}
\begin{itemize}
\item Encoder--decoder loss: $\mathcal{L} = \mathcal{L}_{\mathrm{scVI}} + \lambda\,\mathcal{L}_{\mathrm{adv}}$, with $\mathcal{L}_{\mathrm{scVI}}$ scvi-tools' own loss (ZINB reconstruction plus the KL term at the current KL weight), unchanged; $\lambda$ is the row's lam, constant over the fit \src{\file{scvi_adversarial_plan.py}:361, :373, :436; \file{fit_paper_config.py}:340}.
\item KL weight at 0-based epoch $e$: $e/400$ while $e < 400$ (scvi-tools' default warm-up), so a task with $E < 400$ epochs ends at $(E-1)/400$; X15's completing warm-up has length $E$ and ends at $(E-1)/E$ [S-E3] \src{\file{fit_paper_config.py}:123, :120; SI-44}.
\item Conditioning: conditioned fits register the batch with scVI; unconditioned fits register it in the labels slot, which only the adversary reads \src{\file{scvi_adversarial_plan.py}:534, :537; \file{tests/scvi/test_adversarial_plan_v2.py}:118}.
\end{itemize}
\paragraph{S-M2 Discriminator (JS).}
\begin{itemize}
\item Head: $\ell(z) = W_3\,\sigma(W_2\,\sigma(W_1 z + c_1) + c_2) + c_3 \in \mathbb{R}^{V}$ for $V$ batches, $\sigma$ = ReLU, hidden width 128 \src{\file{wcd/adversarial.py}:50, :76}.
\item Adversary step, one per training step, on the detached input $\bar z$: minimise $\frac{1}{n}\sum_i -\log p_{b_i}(\bar z_i)$ with $p = \mathrm{softmax}(\ell)$ \src{\file{scvi_adversarial_plan.py}:206, :392; \file{wcd/adversarial.py}:71}.
\item Generator term (scvi-tools' fool loss): $\mathcal{L}_{\mathrm{adv}} = -\frac{1}{n}\sum_i \frac{1}{V-1}\sum_{j \neq b_i} \log p_j(z_i)$ \src{\file{scvi_adversarial_plan.py}:292}.
\end{itemize}
\paragraph{S-M3 W1 critics (reference, pooled, barycenter).}
\begin{itemize}
\item Head $C : \mathbb{R}^d \to \mathbb{R}^V$, the same network as S-M2. Estimate of head $k$: $\widehat{W}_k = \frac{1}{n_k}\sum_{i:\,b_i = k} C_k(z_i) - \frac{1}{|T_k|}\sum_{t \in T_k} C_k(t)$; $T_k$ = cells of the reference batch (reference; the reference's own head is inactive), cells of all other batches (pooled), the barycenter support of S-M4 (barycenter) \src{\file{wcd/critic.py}:72, :77, :82}.
\item Critic step, five per training step on the detached input: minimise $-\frac{1}{|A|}\sum_{k \in A}\widehat{W}_k + \mathrm{GP}$, with $A$ the heads whose batch and target are in the minibatch \src{\file{wcd/critic.py}:86--89; \file{scvi_adversarial_plan.py}:394; manifest}.
\item $\mathrm{GP} = 10 \cdot \frac{1}{|A'|}\sum_{k \in A'} \frac{1}{m_k}\sum_{i:\,b_i = k} \bigl(\|\nabla C_k(\hat z_i)\|_2 - 1\bigr)^2$, $\hat z_i = \epsilon_i z_i + (1 - \epsilon_i)\,t_i$, $\epsilon_i \sim U(0, 1)$, $t_i$ drawn uniformly from $T_{b_i}$ \src{\file{wcd/critic.py}:154--214}.
\item Generator term: $\mathcal{L}_{\mathrm{adv}} = \frac{1}{|A|}\sum_{k \in A}\widehat{W}_k$ on the non-detached input, so reference cells move too (a star of pairwise terms) \src{\file{scvi_adversarial_plan.py}:421; \file{scvi_adversarial_plan.py}:10}.
\end{itemize}
\paragraph{S-M4 Barycenter target.}
\begin{itemize}
\item $B^\ast = \arg\min_B \sum_k w_k\, W_2^2(P_k, B)$ with $w_k = 1/K$ over the $K$ batches in the minibatch; $n$ support points with uniform weights, started at the minibatch cells (cold start); at most 10 fixed-point iterations (the update of \citealp{cuturi2014fast}, Algorithm~2 with $\theta = 1$, as implemented in POT free\_support\_barycenter, \citealp{flamary2021pot}; stop threshold $10^{-7}$) with exact transport plans; recomputed from the detached latent at every step \src{SI-14; \file{wcd/barycenter.py}:8--10, :38, :61, :71--73; \file{scvi_adversarial_plan.py}:388}.
\end{itemize}
\paragraph{S-M5 MMD.}
\begin{itemize}
\item $\mathcal{L}_{\mathrm{adv}} = \frac{1}{K}\sum_k \bigl(a_k^\top \mathbf{K} a_k + b_k^\top \mathbf{K} b_k - 2\,a_k^\top \mathbf{K} b_k\bigr)$, with $a_k$ and $b_k$ uniform weights on batch $k$ and on all other cells and $\mathbf{K}_{ij} = \frac{1}{5}\sum_{m \in \{0.25, 0.5, 1, 2, 4\}} \exp\bigl(-\|z_i - z_j\|^2 / (2 d m^2)\bigr)$, $d$ the latent size \src{\file{wcd/alignment.py}:49; \file{wcd/alignment.py}:109--126}.
\item A minibatch with no pair of batches contributes 0, keeping the gradient graph \src{\file{wcd/alignment.py}:69}.
\end{itemize}
\paragraph{S-M6 Sinkhorn divergence.}
\begin{itemize}
\item $\mathcal{L}_{\mathrm{adv}} = \frac{1}{K}\sum_k S_\varepsilon(a_k, b_k)$, $S_\varepsilon(a, b) = \mathrm{OT}_\varepsilon(a, b) - \frac{1}{2}\mathrm{OT}_\varepsilon(a, a) - \frac{1}{2}\mathrm{OT}_\varepsilon(b, b)$, $\mathrm{OT}_\varepsilon$ the regularised transport value $\langle a, f\rangle + \langle b, g\rangle$ from the dual potentials \src{\file{wcd/alignment.py}:232; \file{wcd/alignment.py}:249}.
\item Cost $c(x, y) = \|x - y\|_2$; blur $\varepsilon = 0.05\sqrt{2d}$ (0.224 for $d = 10$), reached by annealing from the cost diameter with ratio 0.8 \src{\file{wcd/alignment.py}:236--243; \file{wcd/alignment.py}:50, :51, :256}.
\item Solve in float64 \src{SI-62}; at most 500,000 iterations at the final $\varepsilon$, with an error at the cap; a NaN marginal error stops at once, $+\infty$ does not \src{SI-61}; float32 was validated against SI-58's criterion and not adopted \src{SI-58; SI-62}; gradient taken once at the converged plan \src{SI-51; \file{wcd/alignment.py}:37}.
\item \pending{convergence tolerance ($10^{-4}$ L1 marginal error, \file{wcd/alignment.py}:56) and whether a fit that keeps reaching the cap counts as failed}{diff review and tag prereg-tier12-v3}
\end{itemize}
\paragraph{S-M7 scvi-tools adversary (scvi\_adv).}
\begin{itemize}
\item scvi-tools' AdversarialTrainingPlan with its classifier adversary (32 hidden units), adversarial loss scaled by $\kappa = 1 -$ KL weight, one classifier step per training step, trained on the posterior sample (scvi-tools' convention, kept for this arm only) \src{\file{scvi_adversarial_plan.py}:5, :177; PREREG \S3}.
\end{itemize}
\paragraph{S-M8 Adversary input and standardisation.}
\begin{itemize}
\item $z^{\mathrm{adv}} = \mu$, the posterior mean (default), or a posterior sample $z$ \src{\file{scvi_adversarial_plan.py}:359; SI-19}.
\item Standardised input (A3; off by default): $(z^{\mathrm{adv}} - \bar\mu)/\max(s_\mu, 10^{-4})$ per dimension, with $\bar\mu$ and $s_\mu$ the minibatch mean and population SD of the posterior means; gradients kept \src{\file{scvi_adversarial_plan.py}:263--270; SI-20}.
\end{itemize}
\paragraph{S-M9 Follow-up arm variants.}
\begin{itemize}
\item R1 penalty (X6; one-vs-rest form of \citealp[Eq.~9]{mescheder2018training}): $\frac{\gamma}{2}\frac{1}{n}\sum_i \|\nabla_z D_{b_i}(z_i)\|^2$, $D_k = \ell_k - \log\sum_{j \neq k} e^{\ell_j}$, $\gamma = 10$, discriminator step only \src{SI-22; \file{wcd/discriminator_losses.py}:30--46}.
\item Spectral normalisation (X6; \citealp{miyato2018spectral}): every head layer spectrally normalised and the gradient penalty dropped \src{\file{wcd/adversarial.py}:50; \file{scvi_adversarial_plan.py}:288}.
\item Reference-anchored JS (X7): binary heads ``batch $k$ against the reference'', class-balanced cross-entropy, non-saturating generator loss with labels flipped, reference cells not detached, one step \src{SI-23; \file{wcd/discriminator_losses.py}:49}. Fixed reference: reference cells detached in the generator step \src{\file{scvi_adversarial_plan.py}:419}. Reference-anchored MMD: MMD$^2$ of each batch to the reference \src{\file{wcd/alignment.py}:129}.
\item Stratified sampler (X8, reference and pooled critics only): $128/V$ cells per batch per minibatch, without replacement within a batch; epoch = $\lceil n_{\mathrm{train}}/128 \rceil$ steps; the fit fails when a batch has fewer training cells than its quota \src{SI-24; \file{wcd/sampling.py}:30}.
\item Importance weights (X3, doses 50, 80 and 95\% only): target = each batch's composition before depletion; $w = 1/(\text{keep fraction})$ for cells of the depleted types in the depleted batch and 1 otherwise; every per-batch mean in the adversary and generator terms becomes the self-normalised mean $\sum_i w_i f(z_i)/\sum_i w_i$, with the same weights in discriminator, reference and pooled critics and MMD \src{SI-25}; the gradient penalty is unweighted \src{\file{wcd/adversarial.py}:91}. \pending{as implemented, $w[b, y] = n_{K_0}(b, y)/n_{K_d}(b, y)$ from the realised nested draws}{the user's confirmation of SI-42 (a lead-agent decision)}
\end{itemize}
\paragraph{S-M10 Rule R1 (A1 to the X1 grids).}
\begin{itemize}
\item Per block (task, decoder, arm), strength $\lambda$ and paired seed $s$: $\Delta b(\lambda)$ = seed mean of $B(\lambda, s) - B(0, s)$, $\Delta c(\lambda)$ likewise; $\sigma_b$ = median over grid points (with at least 2 paired seeds) of the seed SD of $B(\lambda, s) - B(0, s)$, each divided by $\sqrt{\chi^2_{n-1,\,0.5}/(n-1)}$; $\delta_b(\lambda) = \max(0.01,\, 3\sigma_b/\sqrt{n(\lambda)})$; $\sigma_c$, $\delta_c$ likewise \src{PREREG \S2; \file{prereg_rules.py}:393, :405}.
\item No effect: no failed seed, $|\Delta b| < \delta_b$ and $|\Delta c| < \delta_c$. Floor $f$: the largest index with no effect at every point up to it. Collapsed: a failed seed, or $\Delta c \leq -0.10$. Batch peak $p$: the first index above $f$ with $\Delta b \geq \max\Delta b - \delta_b$. Upper end $h$: the smallest index above $f$ from which every point is collapsed or at or after $p$ \src{PREREG \S2}.
\item Edges: when a block reaches an edge, the whole family is extended by the next 2 values of the strength sequence (0.03, 0.01, \ldots{} below; 10000, 30000, \ldots{} above), at most 2 rounds per edge; then the edge is accepted and flagged \src{PREREG \S1, \S2}.
\item Grid: $i_{lo} = \min f$, $i_{hi} = \max h$ over the family's blocks; with at least 6 points, indices $\lfloor i_{lo} + j\,(i_{hi} - i_{lo})/5 + \frac{1}{2} \rfloor$, $j = 0, \ldots, 5$; with fewer, all points plus points added alternately above and below, starting above and skipping a side that has run off the A1 grid \src{PREREG \S2; \file{prereg_rules.py}:448--466}.
\end{itemize}
\paragraph{S-M11 Rules R2 and R3 (A2, A3 to the X1 settings).}
\begin{itemize}
\item Per cell (task, decoder, arm): $D = \mathrm{bio@}b^*(\text{alternative}) - \mathrm{bio@}b^*(\text{default})$, reading seed-mean C at $B = b^*$ by linear interpolation along the cell's window (lo, matched, hi) from rule R4 on A1 \src{PREREG \S3; \file{prereg_rules.py}:610}.
\item Switch only if (1) at least $\lceil 2n/3 \rceil$ of the $n$ cells are evaluable, at least one per task and decoder; (2) the mean $D$ exceeds 0.01; (3) the mean $D$ is positive within each task and decoder; (4) the alternative has no more failed fits. For A2, a sample input that stays below b* in a cell where the mean reaches it keeps the mean \src{PREREG \S3; SI-27; \file{prereg_rules.py}:632}.
\item Cells of a low-signal block stay in $n$ and count as not evaluable \src{SI-56; SI-66}.
\end{itemize}
\paragraph{S-M12 Rule R4 (matched strengths).}
\begin{itemize}
\item Per task and decoder: $b_0$ = seed-mean B of $\lambda = 0$; $M_a$ = arm $a$'s highest failure-free seed-mean B; $b_{\mathrm{common}} = \min_a M_a$ over the six arms; $b^* = b_0 + 0.5\,(b_{\mathrm{common}} - b_0)$; flag bstar\_low\_signal when $b_{\mathrm{common}} - b_0 < 0.02$ \src{SI-21; PREREG \S4; \file{prereg_rules.py}:571, :573}.
\item Matched strength: the failure-free grid point with seed-mean B closest to b*, ties to the smaller $\lambda$; lo and hi are its neighbours. At a grid edge one further sequence value is fitted per round, at most 2 rounds, then flagged edge\_unresolved \src{PREREG \S4}.
\item Variant arms inherit the matched strength of their base arm (discriminator\_sn, discriminator\_r1 and discriminator\_ref from the discriminator; reference\_fixed from the reference critic; pooled\_sn and the pooled critic at other step counts from the pooled critic; mmd\_ref from MMD); subsampled tasks (X3, X8) and architecture variants (X12) inherit the full task's value \src{PREREG \S4; \file{prereg_rules.py}:59}. An inherited strength does not put a variant at matched batch removal \src{D-06}; the follow-up analyses are locked later \src{SI-68}.
\item Rule R1 is frozen before rule R4 on A1 adds any cell-specific A1 point; extension rows of rules R1 and R4 run on their task's machine \src{PREREG \S8; SI-17}.
\end{itemize}
\paragraph{S-M13 P1 reading.}
\begin{itemize}
\item Curve of an arm in a cell: $(b_0, c_0)$ followed by its six registered points $(\bar B_\lambda, \bar C_\lambda)$ in increasing $\lambda$, means over successful runs; a point whose runs all failed is skipped \src{SI-56; SI-69; SI-71}.
\item bio@b*: C at the first point equal to b* (within $10^{-9}$), else linear interpolation in B over the first adjacent pair that brackets b* \src{SI-69; PREREG \S3; \file{prereg_rules.py}:610}.
\item $d = \mathrm{bio@}b^*(\text{reference}) - \mathrm{bio@}b^*(\text{discriminator})$; readings $d_{(i)}$ leave seed $i$ out of both arms and of $\lambda = 0$, with b* fixed; $s^2 = \frac{m-1}{m}\sum_i (d_{(i)} - \bar d_{(\cdot)})^2$ over the $m$ readings that exist \src{SI-69; ANALYSIS\_SPEC App.~A.1, Q2 option A}.
\item Sensitivity analyses: each run read on its own curve \src{SI-69}; points with any failed run removed \src{SI-71}.
\item \pending{not-evaluable rule (low-signal; a curve that never reaches b*; fewer than 4 of 5 leave-one-out readings) and later crossings counted but unused}{diff review (ANALYSIS\_SPEC \S2.2, \S2.3)}
\end{itemize}
\paragraph{S-M14 Trade-off area.}
\begin{itemize}
\item $\mathrm{Area} = \bigl|\bigcup_i [b_0, \bar B_i] \times [c_0 - 0.10, \bar C_i]\bigr|$ over the registered points with $\bar B_i > b_0$ and $\bar C_i > c_0 - 0.10$; pooled as $\mathrm{Area}/(b_{\mathrm{common}} - b_0)$; low-signal cells have no area \src{SI-72; SI-56}.
\item Co-primary with P1, one result per decoder \src{SI-68; SI-70; SI-72 (option A chosen over a secondary area)}.
\item \pending{corner moving with the dropped seed in leave-one-out repeats; corner $(b_0, 0)$ as sensitivity}{diff review (ANALYSIS\_SPEC \S15.3, \S11)}
\end{itemize}
\paragraph{S-M15 Pooling and declarations.}
\begin{itemize}
\item $d_t \sim N(\theta_t, s_t^2)$ with $s_t$ treated as known; $\theta_t = \Delta + u_{f(t)} + v_t$, $u \sim N(0, \tau_f^2)$, $v \sim N(0, \tau_d^2)$; $\Delta \sim N(0, 0.1^2)$, $\tau_f, \tau_d \sim$ half-normal(0.1); computed exactly, without random sampling; prior scales 0.05 and 0.2 as sensitivity \src{SI-73; ANALYSIS\_SPEC \S4}.
\item Direction if $P(\Delta > 0.01) \geq 0.95$ or $P(\Delta < -0.01) \geq 0.95$; practically equivalent if $P(|\Delta| \leq 0.01) \geq 0.95$; otherwise inconclusive \src{SI-74}.
\item Frequentist companion: exact two-sided Wilcoxon signed-rank test on the $d_t$; smallest attainable p with 6 tasks $2/2^6 = 0.03125$, above Holm's first threshold 0.0125 for the four tests P1--P4 \src{PREREG \S5; \file{docs/prereg/r1_record.json} power}; \pending{Holm family of six (0.05/6 = 0.0083) and descriptive use}{diff review (ANALYSIS\_SPEC \S4, \S5, \S15 item 5)}.
\item Power (PREREG \S5 rule; frozen record): the design is flagged UNDERPOWERED; for D = 6 real tasks SE 0.045, minimal detectable effect at 80\% power 0.131 and power at $\Delta = 0.02$ 0.104 (D = 4 families: 0.182 and 0.088); equivalence within $\pm 0.01$ or $\pm 0.02$ is unreachable; computed for the matched-strength contrast on the six A1 blocks, sim1 and the two low-signal blocks included \src{PREREG \S5; \file{docs/prereg/r1_record.json} power}; noise estimates $\hat\sigma = 0.0298$ (12 df) and $\hat\tau = 0.109$ (4 df) \src{\file{docs/prereg/r1_record.json} power}.
\item \pending{integration grid (201 log-spaced points per spread), the zero-spread rule ($s_t < 10^{-6}$), the reported quantities and the pairs-first and six-independent sensitivity analyses}{diff review (ANALYSIS\_SPEC \S4 amendments, \S11, \S15.4)}
\end{itemize}
\paragraph{S-M16 P2 collapse rate.}
\begin{itemize}
\item A run collapses if it failed or if $C \leq C_{\lambda = 0}(\text{same seed}) - 0.10$; rate per arm and decoder over the six registered points and 5 seeds, per task and pooled over the 6 real tasks, with Clopper--Pearson 95\% intervals (\citealp{clopper1934use}) \src{SI-75}.
\item Comparison: per-task rate difference (reference minus discriminator) with the Agresti--Caffo standard error, pooled with the S-M15 model at $\Delta \sim N(0, 0.25^2)$, $\tau \sim$ half-normal(0.25), margin 0.05 \src{SI-75; ANALYSIS\_SPEC App.~A.1, Q9 option A}. \pending{whether SI-75 stands}{diff review}
\end{itemize}
\paragraph{S-M17 Local versus global batch removal.}
\begin{itemize}
\item $L$ = mean of iLISI, kBET and graph connectivity; $G$ = mean of PCR comparison and batch ASW; so $B = (3L + 2G)/5$, and at equal B, $3\Delta L + 2\Delta G = 0$ \src{SI-76; \file{analysis_x1.py}:36, :37}.
\item Question: is $\Delta L > 0$ (critic minus discriminator) at the P1 crossing, with the same interpolation weights; exact two-sided sign test over the real tasks and the S-M15 model, margin 0.01 \src{SI-76}.
\end{itemize}
\paragraph{S-M18 Selection score and selection rules.}
\begin{itemize}
\item $S = 0.4B + 0.6C$, unscaled, for single-strength tables, the equal-tries comparison and the window width \src{SI-77}.
\item \pending{transfer strength: the registered point with the highest mean S over the real tasks of the other families, failure-free, ties to the smaller $\lambda$; best-of-k: expected best of k configurations drawn at random, compared at equal k; window: points with mean S at least the lower 95\% bound of the best point, width in decades}{diff review (ANALYSIS\_SPEC \S8--\S10)}
\end{itemize}

\subsection*{S-D Data}
\paragraph{S-D1 Tasks.} Cells, batches and cell types from the prepared files; keys from each file's record of scIB's keys \src{\file{D/prepped_scib} obs (anndata backed mode); \file{docs/prepped_fingerprints_scib.json}; SI-73 for families}.
{\footnotesize\setlength{\tabcolsep}{3pt}
\begin{longtable}{>{\raggedright\arraybackslash}p{0.15\linewidth}>{\raggedright\arraybackslash}p{0.1\linewidth}rr>{\raggedright\arraybackslash}p{0.14\linewidth}>{\raggedright\arraybackslash}p{0.17\linewidth}r>{\raggedright\arraybackslash}p{0.14\linewidth}}
\toprule
Task & Family & Cells & Batches & Batch key & Label key & Cell types & Modality, organism \\
\midrule
\endhead
\file{pancreas} & pancreas & 16,382 & 9 & tech & celltype & 14 & rna, human \\
\file{lung} & lung & 32,472 & 16 & batch & cell\_type & 17 & rna, human \\
\file{immune} & immune & 33,506 & 10 & batch & final\_annotation & 16 & rna, human \\
\file{immune_hum_mou} & immune & 97,861 & 23 & batch & final\_annotation & 19 & rna, human \\
\file{atac_small} & ATAC & 11,270 & 3 & batchname\_all & final\_cell\_label & 7 & atac, mouse \\
\file{atac_large} & ATAC & 84,813 & 11 & batchname\_all & final\_cell\_label & 7 & atac, mouse \\
\file{sim1} & simulation & 12,097 & 6 & Batch & Group & 7 & rna, none \\
\file{sim2} & simulation & 19,318 & 16 & SubBatch & Group & 4 & rna, none \\
\bottomrule
\end{longtable}
}

\paragraph{S-D2 Provenance and as-distributed counts.} Non-integer batches: batches with any non-integer count, fitted as distributed (SI-15); the Villani batch is the same 1,022 cells in both immune tasks. Genes: in the file; used for fitting. Checksums truncated to 12 hexadecimal characters; full values in \file{docs/prepped_fingerprints_scib.json} \src{\file{docs/prepped_fingerprints_scib.json}; \file{D/prepped_scib} obs}.
{\footnotesize\setlength{\tabcolsep}{3pt}
\begin{longtable}{>{\raggedright\arraybackslash}p{0.15\linewidth}>{\raggedright\arraybackslash}p{0.17\linewidth}>{\raggedright\arraybackslash}p{0.07\linewidth}>{\raggedright\arraybackslash}p{0.11\linewidth}>{\raggedright\arraybackslash}p{0.17\linewidth}ll}
\toprule
Task & Non-integer batches (cells) & Pseudo-time & Genes & Source file & Source MD5 & Counts hash \\
\midrule
\endhead
\file{pancreas} & 4 (5,419): celseq, celseq2, fluidigmc1, smarter & no & 19,093; 2,000 HVG & \file{human_pancreas_norm_complexBatch.h5ad} & \texttt{5aa155652b99} & \texttt{e7f4d76970a0} \\
\file{lung} & 12 (22,771) & no & 15,148; 2,000 HVG & \file{Lung_atlas_public.h5ad} & \texttt{aac0832c3e34} & \texttt{4cee3653e84d} \\
\file{immune} & 1 (1,022): Villani & yes & 12,303; 2,000 HVG & \file{Immune_ALL_human.h5ad} & \texttt{482a6a311743} & \texttt{50def05219df} \\
\file{immune_hum_mou} & 1 (1,022): Villani & yes & 8,135; 2,000 HVG & \file{Immune_ALL_hum_mou.h5ad} & \texttt{3cc1ba9f35b2} & \texttt{cee9b4fcd449} \\
\file{atac_small} & none & no & 3,429; all & \file{small_atac_gene_activity.h5ad} & \texttt{8eb4af0e5e2d} & \texttt{3682994cadf5} \\
\file{atac_large} & none & no & 3,580; all & \file{large_atac_gene_activity.h5ad} & \texttt{bab0bd978406} & \texttt{f6079593c92d} \\
\file{sim1} & none & no & 9,979; 2,000 HVG & \file{sim1_1_norm.h5ad} & \texttt{ea2589604bec} & \texttt{d49b23c37142} \\
\file{sim2} & none & no & 10,000; 2,000 HVG & \file{sim2_norm.h5ad} & \texttt{4d22e971f1e1} & \texttt{e261b5ad83e4} \\
\bottomrule
\end{longtable}
}

\subsection*{S-E Experiments}
\needspace{18\baselineskip}
\paragraph{S-E1 Experiments in the manifest of record.} \file{manifests/paper_manifest_stock_pilot_u5_b10_v3.tsv}, 6,930 rows, SHA-256 6efbe422c8f09f98\ldots; adversary steps per training step in brackets (none: critic-free or no adversary) \src{manifest (pandas count); \file{docs/manifest_v3_counts.csv}}. \pending{X7 and X8 rows}{X7/X8 redesign (SI-68)}
{\footnotesize\setlength{\tabcolsep}{3pt}
\begin{longtable}{>{\raggedright\arraybackslash}p{0.05\linewidth}r>{\raggedright\arraybackslash}p{0.18\linewidth}>{\raggedright\arraybackslash}p{0.24\linewidth}>{\raggedright\arraybackslash}p{0.11\linewidth}>{\raggedright\arraybackslash}p{0.1\linewidth}>{\raggedright\arraybackslash}p{0.16\linewidth}}
\toprule
Exp. & Rows & Tasks & Arms [steps] & Decoders & Seeds & Strength source \\
\midrule
\endhead
A1 & 1,098 & atac\_small, immune, sim1 & barycenter\,[5], discriminator\,[1], mmd, none, pooled\,[5], reference\,[5], sinkhorn & both & 100--102 & A1 grid 0.1--3000 (10 values) \\
A2 & 108 & atac\_small, immune & discriminator\,[1], pooled\,[5], reference\,[5] & both & 100--102 & rule R4 on A1: lo, matched, hi \\
A3 & 180 & atac\_small, immune & discriminator\,[1], mmd, pooled\,[5], reference\,[5], sinkhorn & both & 100--102 & rule R4 on A1: lo, matched, hi \\
X1 & 3,000 & all 8 & barycenter\,[5], discriminator\,[1], mmd, none, pooled\,[5], reference\,[5], scvi\_adv, sinkhorn & both & 0--4 & rule R1 grid g1--g6 [S-E4]; 0 for none, scvi\_adv \\
X13 & 408 & all 8 & harmony, pca, scanorama, scanvi, sysvi & cond. (scANVI, sysVI); none (CPU) & 0--4 (scANVI, sysVI); 0 (CPU) & method knob ($\theta$, knn, cycle weight) \\
X3 & 828 & atac\_small, immune\_hum\_mou, pancreas, sim2 & discriminator\,[1], mmd, none, pooled\,[5], reference\,[5] & cond. & 10--12 & rule R4 on X1: lo, hi; 0 for none \\
X6 & 216 & atac\_small, immune, pancreas & discriminator\,[1], discriminator\_r1\,[1], discriminator\_sn\,[1], pooled\,[1/5/10], pooled\_sn\,[1/5] & cond. & 10--12 & rule R4 on X1: lo, matched, hi \\
X7 & 216 & atac\_small, pancreas, sim1 & discriminator\_ref\,[1], mmd\_ref, reference\,[5], reference\_fixed\,[5] & cond. & 10--12 & rule R4 on X1: matched \\
X8 & 504 & immune\_hum\_mou, lung, sim2 & discriminator\,[1], mmd, none, pooled\,[5], reference\,[5] & cond. & 10--12 & rule R4 on X1: matched; 0 for none \\
X12 & 288 & atac\_small, immune, pancreas & discriminator\,[1], reference\,[5] & cond. & 10--12 & rule R4 on X1: lo, hi \\
X15 & 84 & atac\_large, immune & discriminator\,[1], mmd, none, pooled\,[5] & cond. & 10--12 & rule R4 on X1: lo, hi; 0 for none \\
Total & 6,930 & & & & & \\
\bottomrule
\end{longtable}
}
\begin{itemize}
\item X3 targets, cells per dose: atac\_small, Fang et al. - CEMBA180305\_2B, excitatory and inhibitory neurons, 8,402 \src{SI-33; SI-41}; pancreas, inDrop3, alpha and acinar cells, 14,409 \src{SI-41; manifest}; immune\_hum\_mou, MCA\_BM\_2, neutrophils and monocyte progenitors, 20,000 \src{SI-41; manifest}; sim2, Batch3Sub1, Group1, 17,872 \src{SI-32; SI-41}.
\end{itemize}

\paragraph{S-E2 Settings.}
{\small
\begin{longtable}{>{\raggedright\arraybackslash}p{0.24\linewidth}>{\raggedright\arraybackslash}p{0.42\linewidth}>{\raggedright\arraybackslash}p{0.26\linewidth}}
\toprule
Setting & Value & Source \\
\midrule
\endhead
Latent size & 10 (X12: 10 or 30) & \src{SI-10; manifest} \\
Hidden layers, units & 1, 128 & \src{SI-10; manifest} \\
Likelihood & ZINB (sysVI: Gaussian, on the file's normalised values) & \src{SI-10; \file{fit_paper_config.py}:404--409} \\
Training split & 0.9 & \src{SI-10; manifest} \\
Minibatch & 128 (X12: 128 or 512) & \src{SI-10; manifest} \\
Epochs & $\min(400, \mathrm{round}(20000 \times 400/n))$ [S-E3] & \src{SI-04; \file{build_paper_manifest.py}:114} \\
Early stopping & off & \src{\file{scvi_adversarial_plan.py}:563} \\
Generator optimiser & Adam, $10^{-3}$, $\epsilon$ 0.01, weight decay $10^{-6}$ & \src{\file{scvi_adversarial_plan.py}:245} \\
KL warm-up & 400 epochs, linear (X15: $E$ epochs) & \src{\file{fit_paper_config.py}:123; SI-44} \\
Adversary head & 128, 128, ReLU (X12 width: 32 or 128) & \src{\file{wcd/adversarial.py}:50; manifest} \\
Discriminator optimiser, steps & Adam, $10^{-3}$, $\epsilon$ 0.01, weight decay $10^{-6}$; 1 & \src{\file{scvi_adversarial_plan.py}:260; PREREG \S9, 2026-10-04 (CR-12); SI-01; manifest} \\
Critic optimiser, steps & Adam, $10^{-4}$, betas (0, 0.9); 5 (X6 pooled: 1, 5, 10) & \src{\file{scvi_adversarial_plan.py}:258; manifest} \\
Adversary learning rate (X12) & $10^{-4}$ or $10^{-3}$ & \src{manifest} \\
Gradient penalty & weight 10, two-sided & \src{\file{wcd/critic.py}:158--210} \\
Barycenter & at most 10 cold-start iterations (stop threshold $10^{-7}$), support = minibatch & \src{SI-14; \file{wcd/barycenter.py}:38, :72--73} \\
MMD kernel & Gaussian, bandwidths $\sqrt{d} \times \{0.25, 0.5, 1, 2, 4\}$ & \src{\file{wcd/alignment.py}:49} \\
Sinkhorn & blur $0.05\sqrt{2d}$, annealing ratio 0.8, float64, cap 500,000 & \src{\file{wcd/alignment.py}:50--51; SI-61; SI-62} \\
R1 penalty (X6) & $\gamma = 10$ & \src{SI-22; manifest} \\
Adversary input & posterior mean (rules R2/R3 cannot switch while immune, conditioned, is flagged low-signal; methods); \pending{formal rule R2 re-run}{after the Sinkhorn redo} & \src{SI-19; SI-66; \file{docs/prereg/r4_a1_record.json}; manifest} \\
Standardisation & off (rules R2/R3 cannot switch while immune, conditioned, is flagged low-signal; methods); \pending{formal rule R3 re-run}{after the Sinkhorn redo} & \src{SI-20; SI-66; \file{docs/prereg/r4_a1_record.json}; manifest} \\
Reference batch & highest cell-type entropy (X3: dose-0 choice; X7: every batch) & \src{\file{wcd/data.py}:32; SI-31; manifest} \\
scANVI epochs & $\min(10, \max(2, \mathrm{round}(E/3)))$ after scVI & \src{\file{fit_paper_config.py}:401} \\
sysVI cycle weight & 1, 2, 5, 10, 20, 50 & \src{\file{build_paper_manifest.py}:197; manifest} \\
Harmony $\theta$, PCs & 0, 0.5, 1, 2, 4, 8 at 10; 2 at 50 & \src{SI-34} \\
Scanorama knn, dimensions & 2, 5, 10, 20, 40, 80 at 10; 20 at 100 & \src{SI-46} \\
PCA components & 10; 50 & \src{SI-36} \\
Scoring graph & scanpy neighbours on the embedding (default 15 neighbours) & \src{\file{score_scib_native.py}:207; scanpy 1.11.5 neighbors()} \\
kBET seed & 0, set in R before each call & \src{\file{score_scib_native.py}:54, :80} \\
\bottomrule
\end{longtable}
}

\paragraph{S-E3 Epochs and final KL weight per task.} Epochs from the manifest (equal to the heuristic for every task); final KL weight = KL weight in the last epoch, $(E-1)/400$ for $E \leq 400$; completing warm-up only in X15 \src{manifest; \file{fit_paper_config.py}:123; SI-44}.
{\small
\begin{longtable}{lrrrr}
\toprule
Task & Cells & Epochs $E$ & Final KL weight (stock) & Final KL weight (X15 completing) \\
\midrule
\endhead
\file{pancreas} & 16,382 & 400 & 0.998 & -- \\
\file{lung} & 32,472 & 246 & 0.613 & -- \\
\file{immune} & 33,506 & 239 & 0.595 & 0.996 \\
\file{immune_hum_mou} & 97,861 & 82 & 0.203 & -- \\
\file{atac_small} & 11,270 & 400 & 0.998 & -- \\
\file{atac_large} & 84,813 & 94 & 0.233 & 0.989 \\
\file{sim1} & 12,097 & 400 & 0.998 & -- \\
\file{sim2} & 19,318 & 400 & 0.998 & -- \\
\bottomrule
\end{longtable}
}

\paragraph{S-E4 Frozen X1 strength grids (rule R1).} Six strengths per family and decoder from \file{docs/prereg/r1_record.json} (status frozen, git commit c88cce3) \src{\file{docs/prereg/r1_record.json} families[*].x1\_grid; SI-18}. The record still holds a Sinkhorn grid that is superseded \src{PREREG \S9, 2026-10-13 (Sinkhorn entry); SI-57}.
{\small
\begin{longtable}{>{\raggedright\arraybackslash}p{0.3\linewidth}>{\raggedright\arraybackslash}p{0.14\linewidth}>{\raggedright\arraybackslash}p{0.3\linewidth}>{\raggedright\arraybackslash}p{0.16\linewidth}}
\toprule
Family (arms) & Decoder & Six registered strengths & Source \\
\midrule
\endhead
JS (discriminator) & conditioned & 0.1, 0.3, 1, 3, 10, 30 & \src{\file{r1_record.json} js$|$1} \\
JS (discriminator) & unconditioned & 0.03, 0.1, 0.3, 1, 3, 10 & \src{\file{r1_record.json} js$|$0} \\
W1 (reference, pooled, barycenter) & conditioned & 0.1, 0.3, 3, 10, 100, 300 & \src{\file{r1_record.json} w1$|$1} \\
W1 (reference, pooled, barycenter) & unconditioned & 0.3, 1, 3, 10, 30, 100 & \src{\file{r1_record.json} w1$|$0} \\
MMD (mmd) & conditioned & 3, 10, 30, 300, 1000, 3000 & \src{\file{r1_record.json} mmd$|$1} \\
MMD (mmd) & unconditioned & 3, 10, 30, 100, 300, 1000 & \src{\file{r1_record.json} mmd$|$0} \\
Sinkhorn (sinkhorn) & conditioned & \pending{six strengths}{rule R1 re-run on the refitted Sinkhorn fits (SI-57)} & -- \\
Sinkhorn (sinkhorn) & unconditioned & \pending{six strengths}{rule R1 re-run on the refitted Sinkhorn fits (SI-57)} & -- \\
\bottomrule
\end{longtable}
}
The unconditioned JS grid starts at 0.03, a value added to A1 by one round of rule R1's low-edge extension \src{\file{r1_record.json} js$|$0 rounds, a1\_grid}.

\paragraph{S-E5 Software.} Package versions of the fitting and scoring environments; the committed \file{environment.yaml} and \file{conda-lock.yml} describe an older general environment, not these \src{git log 08f1e3c -- \file{environment.yaml} \file{conda-lock.yml}: last change 2025-06-21}. The JHPCE fitting environment was rebuilt with conda packages and pip versions equal to the local one, apart from 34 declared CUDA builds \src{SI-48; NB-20261004-20}; each latent records its scvi-tools and PyTorch versions and each score row its scorer versions \src{\file{fit_paper_config.py}:353--355; \file{score_scib_native.py}:38--41}. Versions agree with the provenance written into all 1,098 A1 latents and score rows \src{\file{D/tier12_v2/A1} latents and scores}.
{\small
\begin{longtable}{>{\raggedright\arraybackslash}p{0.2\linewidth}>{\raggedright\arraybackslash}p{0.18\linewidth}>{\raggedright\arraybackslash}p{0.26\linewidth}>{\raggedright\arraybackslash}p{0.24\linewidth}}
\toprule
Environment & Package & Version & Source \\
\midrule
\endhead
Fitting (env scvi-api) & Python & 3.11.15 & \src{\file{docs/jhpce/gate_data/versions_fit_local.json}} \\
scvi-api & scvi-tools & 1.4.2 & \src{\file{cluster/envspec/local_scvi-}\file{api_pip_freeze.txt}} \\
scvi-api & PyTorch & 2.13.0 (CUDA 13.0 build locally, 12.6 on L40S) & \src{pip freeze; \file{versions_fit_local.json}; \file{gate_report_L40S/versions_fit_gpu.json}} \\
scvi-api & Lightning & 2.6.5 & \src{pip freeze} \\
scvi-api & POT & 0.9.7.post1 & \src{pip freeze} \\
scvi-api & anndata, scanpy, NumPy, SciPy & 0.12.19, 1.11.5, 2.4.6, 1.17.1 & \src{pip freeze} \\
Scoring and CPU baselines (env wcd-kbet) & Python & 3.11.15 & \src{\file{docs/jhpce/gate_data/versions_score_local.json}} \\
wcd-kbet & scib & 1.1.7 & \src{\file{cluster/envspec/local_wcd-}\file{kbet_pip_freeze.txt}} \\
wcd-kbet & R, kBET & 4.5.2 (2025-10-31), 0.99.6 & \src{\file{versions_score_local.json}} \\
wcd-kbet & scanpy, anndata, NumPy & 1.11.5, 0.11.4, 1.26.4 & \src{pip freeze} \\
wcd-kbet & Scanorama & 1.7.4 & \src{pip freeze} \\
wcd-kbet & Harmony (harmony-pytorch) & 0.1.7 in env wcd-kbet; missing from the committed freeze, which lists harmonypy 2.0.0 (not used) & \src{\file{run_cpu_baselines.py}:76, :142 (version written into each latent); \file{cluster/envspec/local_wcd-}\file{kbet_pip_freeze.txt}:34} \\
\bottomrule
\end{longtable}
}

\paragraph{S-E6 Hardware placement, fit and score provenance.}
{\small
\begin{longtable}{>{\raggedright\arraybackslash}p{0.3\linewidth}>{\raggedright\arraybackslash}p{0.36\linewidth}>{\raggedright\arraybackslash}p{0.26\linewidth}}
\toprule
Tasks & Machine & Source \\
\midrule
\endhead
atac\_small, immune, sim1 & local NVIDIA GeForce RTX 3080, 10 GiB & \src{SI-30; \file{versions_fit_local.json}} \\
atac\_large, immune\_hum\_mou, lung, pancreas, sim2 & JHPCE, NVIDIA L40S (2 GPUs) & \src{SI-39; \file{gate_report_L40S/versions_fit_gpu.json}} \\
all (scoring, CPU baselines) & local CPU & \src{SI-47; PREREG \S9, 2026-10-04 (b)} \\
\bottomrule
\end{longtable}
}
\begin{itemize}
\item One task, one machine: every arm, seed and strength of a task is fitted on the same machine \src{SI-17}.
\item Limits of the commit-equivalence check: score rows carry the scorer's commit, not the fit's, and Sinkhorn rows are refused only at c88cce3, so each stage is fitted and scored from one worktree pinned at the tag \src{\file{docs/EQUIVALENT_COMMITS.json} meaning; PREREG \S9, 2026-10-04 (b)}.
\item scib's silent fallbacks (trajectory root cell set to 0; kBET labels forced to 1 or skipped) are flagged in every score row and their counts per rule cell are reported with the results; no rule reads them \src{PREREG \S9, 2026-10-04 (CR-05, CR-08)}.
\item 250 rule-free fits of the five JHPCE tasks (X1 none and scvi\_adv; X13 scANVI and sysVI) were fitted at c88cce3 in two jobs of one L40S each and harvested for local scoring \src{SI-47; NB-20261006-02}; they are non-Sinkhorn rows, covered by the commit equivalence \src{SI-54}.
\end{itemize}

\subsection*{S-X Disclosures}
\begin{enumerate}
\item A code check of tag prereg-tier12-v1 led to fixes before tag prereg-tier12-v2 \src{SI-40; PREREG \S9, 2026-10-04}. A1 restarted from zero at tag prereg-tier12-v2; the 202 A1 fits made at prereg-tier12-v1 stay on disk, unused, and no A1 metric value from them informed any decision \src{SI-43}.
\item The strength ranges of X1 come from 3 of the 8 tasks (atac\_small, immune, sim1), fitted on seeds 100--102, disjoint from X1; no A1 fit is used as an X1 fit \src{PREREG, Disclosure (for Methods); SI-13}.
\item Strength-grid edges: X1 tasks whose matched strength falls on a grid edge are extended by rule R4 for matching only, and edge optima are disclosed per task \src{PREREG \S8; SI-55}.
\item Pairing: equal seeds give equal initialisation and training split in every arm, but the minibatch order is shared only within an arm and with the arms without an adversary head (MMD, Sinkhorn) \src{PREREG \S9, 2026-10-04 (CR-07)}. Declared differences between arms (critic and discriminator steps; scANVI's extra epochs; sysVI without KL warm-up) are stated in Methods \src{PREREG \S9, 2026-10-04 (CR-12)}.
\item The Scanorama grid was amended before any X13 fit (knn 160 replaced by 2 after knn 160 did not finish within 8 h on immune\_hum\_mou); the X13 grid was not chosen from results \src{SI-37; SI-46}.
\item If rules R2 and R3 both switch, the untested combination is disclosed \src{PREREG \S3}.
\item Rule R1 depends on the trajectory metric through its noise estimate; the registered R1 and frozen grids are kept and R1 without trajectory is reported as a sensitivity analysis \src{SI-53}. \pending{result of R1 without trajectory}{its Sinkhorn part recomputed after the Sinkhorn redo}
\item In the Sinkhorn arm at tag prereg-tier12-v2 the gradient was counted twice and the solve stopped before converging at the final smoothing step \src{SI-51; I031; I032}.
\item The 234 Sinkhorn pilot fits made with it (A1 180, R1 extension 18, A3 36) are replaced by refits in rule order, in which the numbers of extension and A3 rows are outputs of rules R1 and R4 on the new fits \src{SI-51; SI-57}.
\item Part of X1 was fitted and scored before the X1 analyses were specified: the local run of 2026-10-12 stopped with 19 X1 rows scored (17 adversarial-arm rows, all immune, conditioned decoder, seed 0, plus none and scvi\_adv) \src{NB-20261012-05; NB-20261014-21}.
\item An audit of the project's chat records found no display or analysis of these values \src{NB-20261014-21}. \pending{the user's confirmation that nobody read an X1 adversarial-arm score, and the definition of ``examined''}{diff review (ANALYSIS\_SPEC App.~A.3, item 2)}
\item The area corner (SI-72), the prior scales (SI-73) and the P2 margin and prior (SI-75) were chosen after the A1 pilot ranges and collapse rates had been shown with the questions \src{NB-20261014-05; NB-20261014-07; NB-20261014-09; ANALYSIS\_SPEC App.~A.1}.
\item The ledger decisions SI-51 to SI-77 (2026-10-13 and 2026-10-14) were made after A1, A2 and A3 had been fitted and scored and rules R2 and R3 first run (NB-20261012-01), and after the local X1 run began (NB-20261012-04) \src{\file{CONSTRAINTS.md} dates; NB-20261012-01; NB-20261012-04}.
\item sysVI rows carry likelihood and decoder values that the sysVI fit does not read; per-arm refusals of unread settings are kept as implemented \src{SI-63}.
\item Final pre-publication review (2026-10-13): 99 unique issues (65 confirmed, 32 partly confirmed, 2 not reproduced), 17 rated fix-before-launch; its fixes are part of version 3 \src{\file{docs/review_final/REVIEW_FINAL.md}}. \pending{the fixes as frozen}{diff review and tag prereg-tier12-v3 (SI-52)}
\end{enumerate}


\end{document}
